\providecommand{\toplevelprefix}{../..}  %
\documentclass[../../book-main_zh.tex]{subfiles}

\begin{document}

\chapter{作为展开优化的深度表示}
\label{ch:deep-networks}

\begin{quote}
    \hfill    ``{\em 凡我不能创造的，我就不能理解}。''

    $~$ \hfill --- 理查德·费曼 (Richard Feynman)
\end{quote}
\vspace{5mm}

在前面的章节中，我们展示了如何识别高维空间中的低维结构，\textit{主要侧重于线性结构}。
例如，当观测数据 $\bm{x}$ 服从统计模型 $\bm{x} = \bm{U}\bm{z} + \bm{\varepsilon}$ 时，我们引入了主成分分析 (PCA) 来学习线性去噪器 $\hat{\bm{U}}^{\top}$。
在这种设定下，学习到的表示是经过线性变换的输入数据 $\hat{\bm{U}}^{\top}\bm{x}$。
在线性模型假设下，人们可以通过高效的优化算法和强有力的理论保证来学习低维线性结构。
此外，线性模型假设涵盖了广泛的应用和问题，包括人脸识别、磁共振图像恢复以及结构纹理恢复~\cite{Wright-Ma-2022}。

另一方面，在处理实际应用时，尤其是当输入数据 $\bm{x}$ 非常复杂（如语音和自然语言、图像和视频以及机器人运动）时，线性模型可能会受到限制。这类数据的低维分布通常是{\em 非线性的}。
如何处理非线性在控制理论、信号处理和模式识别等不同学科中有着悠久的历史。人们付出了大量努力试图将线性模型的方法和解决方案扩展到处理非线性问题，包括早期将 PCA 扩展到非线性 PCA 的尝试（正如我们将在\Cref{ch:consistent-representations}中详细研究的那样）。在大多数情况下，这些方法是基于对数据分布的特定假设而设计的，并针对特定问题进行了量身定制。

最近，深度神经网络在广泛的数据和应用中取得了显著的成功。一个神经网络
\begin{equation}
    f(\cdot, \bm \theta) \colon \x
    \xrightarrow{\hspace{1mm} f^{\mathrm{pre}} \hspace{1mm}} \z^1
    \rightarrow \cdots
    \rightarrow \z^\ell \xrightarrow{\hspace{1mm} f^\ell \hspace{1mm}}
    \z^{\ell+1} \rightarrow  \cdots \to \z^{L+1} = \z.
\end{equation}
能够为下游应用学习有效的特征/表示。例如，训练好的深度神经网络 $f(\cdot, \bm \theta)$ 可用于将图像映射为特征向量，即 $\bm{z}_i = f(\bm{x}_i,\bm \theta)$，同时可以在这些表示 $\{\bm{z}_i\}$ 之上学习一个线性分类器。一个著名的突破是 AlexNet~\cite{krizhevsky2012imagenet}，这是一个使用超过一百万张自然图像训练的深度卷积神经网络，其性能超越了以往所有基于手工特征的方法。AlexNet 与以往方法的一个关键区别在于，前者\textit{从海量数据中学习非线性变换的参数}，并通过反向传播 (BP) \cite{Back-Prop} 进行训练，详见\Cref{app:optimization}的\Cref{app:BP-section}。

随后的流行做法是使用其他凭经验设计的人工深度神经网络对映射 $f$ 进行建模，并通过 BP 从随机初始化中学习参数 $\bm \theta$。从 AlexNet \cite{krizhevsky2012imagenet} 开始，现代深度网络的架构不断地被凭经验进行修改和改进。诸如 VGG \cite{simonyan2014very}、ResNet \cite{he2016deep}、DenseNet \cite{dense-net}、CNN、RNN 或 LSTM \cite{LSTM}、Transformer \cite{vaswani2017attention} 以及混合专家模型 (MoE) \cite{MoE,Fedus-2022} 等网络架构不断突破性能极限。作为提高深度网络性能努力的一部分，网络的几乎每个组件都经过了经验性的审查，并提出了各种修改和改进。它们不局限于非线性激活函数 \cite{maas2013rectifier,klambauer2017self,xu2015empirical,nwankpa2018activation}、跳跃连接 \cite{ronneberger2015u,he2016deep}、归一化 \cite{ioffe2015batch,ba2016layer,ulyanov2016instance,wu2018group,miyato2018spectral}、上/下采样或池化 \cite{scherer2010evaluation}、卷积~\cite{lecun1998gradient,krizhevsky2012imagenet} 等。然而，几乎所有这些修改都是通过多年的经验性{试错}或消融实验发展而来的。最近的一些实践甚至走向了极端，通过广泛的随机搜索技术来寻找有效的网络结构和训练策略，例如神经架构搜索 (NAS) \cite{NAS-1,Baker2017DesigningNN}、AutoML \cite{automl} 和元学习 (Learning to Learn) \cite{andrychowicz2016learning}。

尽管深度神经网络得到了广泛应用，但目前尚不清楚这种构建网络背后的基本设计原则是什么。特别是，尚不清楚网络的每一层执行的是什么数学功能。在本章中，基于前几章的结果，我们将开发一个原则性的框架，该框架将为深度网络（包括其各个层以及整个网络）的作用提供完全严格的数学解释。

为了理解深度网络以及如何更好地设计它们，我们必须从表征学习的目标开始。在前几章中，我们已经论证了该目标是识别本质上低维的数据分布，然后将其转换为紧凑且结构化（例如分段线性）的表示。正如我们在上一章中所看到的，识别低维数据分布的一般方法是通过一个逐步最小化分布的熵或编码率的压缩过程。然而，到目前为止，我们一直使用凭经验设计的深度网络来建模或近似旨在优化这些目标的操作，例如用于去噪的得分函数（在\Cref{sec:denoising-intro}中）或最大化编码率衰减的变换（在\Cref{sec:MCR-landscape}中）。

正如我们在上一章（特别是\Cref{sec:chap4-representation-learning-problem}）中所论证的那样，人们可以通过从“惰性 (lazy)”表示（将所有数据建模为一个大高斯分布）中获得的表示信息来衡量所得表示的优劣。\footnote{我们在上一章中已经看到，这是对采样数据集解释的一种特定选择。}特别是，如果\textit{我们使用高斯混合（子空间）\footnote{我们在上一章中已经对其进行了深入研究。}作为原型分布来近似感兴趣的非线性分布}，那么我们可以使用相关高斯分布的率失真函数（之和）来有效地测量这种表示的编码率。那么，通过这种建模所获得的{\em 信息增益}或减少的（相对）熵，可以通过惰性表示的编码率与更精细表示的编码率之间的差值来衡量。因此，表征学习的目标就是最大化这种信息增益，这也称为编码率衰减目标。

正如我们将在本章中看到的，一旦表征学习的目标明确，深度神经网络的作用恰恰就是帮助迭代地优化该目标。深度神经网络的每一层都可以自然地推导为增量最大化信息增益的迭代优化步骤，这包括流行的 ResNet、CNN 和 Transformer 架构，以及其他更高级的变体。特别是，本章旨在回答关于深度网络的以下问题：
\begin{itemize}
    \item \Cref{sec:chap4-white-box-model-via-unrolling} --- 给定学习表示的优劣度量，如何通过针对该目标的展开优化来构建从数据到最优表示的非线性映射？
    \item \Cref{sec:chap4-white-box-transformer} --- 上述展开方法将如何为流行的 Transformer 架构提供原则性的解释；如果是这样，相关的目标和优化机制是什么？
    \item \Cref{sec:chap4-derive-white-box-transformer-variants} --- 该框架将如何指导我们设计更高效或更精简的深度架构？
\end{itemize}

\section{通过展开优化构建白盒深度网络}\label{sec:chap4-white-box-model-via-unrolling}

现在，如果我们同意如\Cref{sec:chap4-representation-learning-problem}中所讨论的，最大化编码率衰减或信息增益会产生所需的表示，那么剩下的问题就是如何构建和学习从数据 $\X$ 到最优表示 $\Z^*$ 的（非线性）映射。这涉及设计一种网络架构和学习算法，使其能够有效地捕获数据中的潜在结构并忠实地实现最优表示。

\subsection{源自展开梯度下降的深度网络}

在上一章中，我们提出了编码率衰减目标 \eqref{eqn:maximal-rate-reduction}，作为学习数据线性判别表示的原则性目标。然而，我们尚未指定用于从输入数据 $\x$ 中提取此类表示的特征映射 $\z = f(\x, \bm \theta)$ 的架构。一个直接的选择是使用传统的深度网络（如 ResNet）来实现 $f(\x, \bm \theta)$。正如我们在\Cref{eg:Rate-Reduction-CIFAR10}中所看到的，这种选择在经验上通常会带来不错的性能。尽管如此，采用任意深度网络仍然存在几个未解决的问题。虽然学习到的特征表示现在更具可解释性，但网络本身仍然{\em 不具备}可解释性。目前尚不清楚为什么任何选定的“黑盒”网络都能够优化所需的 MCR$^2$ 目标。良好的经验结果（例如使用 ResNet）并不一定能证明网络架构和算子的特定选择是合理的：为什么甚至需要深度分层模型；额外的层试图改进或简化什么；多宽和多深才足够；或者对于所使用的卷积（以流行的多通道形式）和非线性算子（例如 ReLU 或 softmax）是否有任何严格的理论依据？

在本章中，我们将展示使用梯度上升来最大化 \eqref{eqn:maximal-rate-reduction} 中定义的编码率衰减 $\Delta R_{\epsilon}(\Z \mid \bm \Pi)$，会自然地导出一个实现所需映射的“白盒”深度网络。所有网络层、线性/非线性算子和参数都是{\em 以纯前向传播的方式显式构建的}。此外，这种网络架构类似于现有的凭经验设计的深度网络，为其设计提供了原则性的理论依据。

\paragraph{用于编码率衰减的梯度上升。} 从上一章中我们可以看到，为了寻找线性判别表示 (LDR)，在数学上，我们本质上是在寻找一个连续映射 $f(\cdot): \x \mapsto \z$，将数据 $\X = [\x_1, \ldots, \x_N] \in \Re^{D \times N}$（或从数据中提取的初始特征\footnote{正如我们将在下一节中看到这种特征提取的必要性。}）映射到最优表示 $\Z = [\z_1, \ldots, \z_N] \in \Re^{d \times N}$，从而最大化以下编码率衰减目标（另见 \eqref{eq:MCR pi}）：
\begin{equation}\label{eq:mcr-formulation}
    \begin{split}
        \Delta R_{\epsilon}(\Z \mid \bm{\Pi}) \doteq
        \underbrace{\frac{1}{2}\log\det \Big(\I + {\alpha} \Z
        \Z^{\top} \Big)}_{R_{\epsilon}(\Z)} \;-\;
        \underbrace{\sum_{k=1}^{K}\frac{\gamma_{k}}{2}
            \log\det\Big(\I + {\alpha_{k}} \Z \bm{\Pi}_{k} \Z^{\top}
        \Big)}_{R_{\epsilon}^c(\Z \mid \bm \Pi)},
    \end{split}
\end{equation}
其中 $\epsilon > 0$ 是规定的量化误差，为简便起见，我们记作\footnote{请注意，与\Cref{ch:denoising}相比，我们使用了稍微简化的符号。}
\begin{align*}
    \alpha \doteq \frac{d}{N\epsilon^2}, \quad \alpha_{k} \doteq
    \frac{d}{\mathrm{tr}(\bm{\Pi}_{k})\epsilon^2}, \quad \gamma_{k}
    \doteq \frac{\mathrm{tr}(\bm{\Pi}_{k})}{N}, \quad \text{for}\ k
    = 1,\ldots, K.
\end{align*}

问题实际上归结为是否存在一种{\em 构造性}的方法来找到这种从 $\bm x$ 到 $\bm z$ 的连续映射 $f(\cdot,\bm \theta)$？为此，让我们考虑增量地最大化作为 $\Z \subseteq \mathbb{S}^{d-1}$ 函数的目标 $\Delta R_{\epsilon}$。尽管可能有许多优化方案可供选择，但为简单起见，我们首先考虑可以说是最简单的投影{\em 梯度上升} (PGA) 方案：\footnote{请注意，我们在 $\Z_j$ 上使用下标 $j$ 来表示第 $j$ 类的特征，在 $\Z^{\ell}$ 上使用上标 $\ell$ 来表示第 $\ell$ 次迭代或第 $\ell$ 层的所有特征。}
\begin{equation}
    \bm Z^{\ell+1}   \; \propto \; \bm Z^{\ell} + \eta \cdot
    \frac{\partial \Delta R_{\epsilon}}{\partial \bm Z}(\Z^{\ell})
    \quad \mbox{s.t.} \quad \Z^{\ell+1} \subseteq \mathbb{S}^{d-1},
    \quad \ell = 1, 2, \ldots,
    \label{eqn:gradient-descent}
\end{equation}
对于某个步长 $\eta >0$，迭代从给定数据 $\bm Z^{0} = \bm X$ 开始。\footnote{同样，为简单起见，我们在这里首先假设初始特征 $\bm Z^{1}$ 就是数据本身。请注意，这里的 $\ell$ 表示迭代次数。因此，数据和特征具有相同的维度 $d$。但这并非必须如此。正如我们将在下一节中看到的，初始特征可以是从数据的一些（提升的）特征开始，并且原则上可以具有不同（高得多）的维度。所有后续迭代具有相同的维度。}
该方案可以解释为应如何增量地调整当前特征 $\Z^\ell$（初始化为输入数据 $\bm X$）的位置，以便生成的 $\Z^{\ell +1}$ 能够改善编码率衰减 $\Delta R_{\epsilon}$，如\Cref{fig:gradient-flow}所示。
\begin{figure}
    \centering
    \includegraphics[width=0.85\linewidth]{\toplevelprefix/chapters/chapter5/figs/redu_gradient_diagram.png}
    \caption{通过梯度流进行增量形变，既将每个类的数据展平到一个子空间中，又将不同的类推开。}
    \label{fig:gradient-flow}
\end{figure}

简单的计算表明，梯度 ${\partial \Delta
R_{\epsilon}}/{\partial \bm Z}$ 需要计算 $\Delta R_{\epsilon}$ 中两项的以下导数：
\begin{equation}\label{eqn:expand-directions}
    \frac{1}{2}\frac{\partial \log \det (\I \!+\! \alpha \Z \Z^{\top}
    )}{\partial \bm Z}(\Z^\ell) = \underbrace{\alpha(\I \!+\!
    \alpha\Z^\ell (\Z^\ell)^\top)^{-1}}_{\bm E^{\ell} \; \in
    \Re^{d\times d}}\Z^\ell,
\end{equation}

\begin{equation}\label{eqn:compress-directions}
    \frac{1}{2}\frac{\partial \left( \gamma_{k}  \log \det (\I +
    \alpha_{k} \Z \bm \Pi_k \Z^{\top} )  \right)}{\partial \bm
    Z}(\Z^\ell) = \gamma_{k}  \underbrace{ \alpha_{k}  (\I +
    \alpha_{k} \Z^\ell \bm \Pi_k (\Z^\ell)^{\top})^{-1}}_{\bm
    C^{\ell}_k \; \in \Re^{d\times d}} \Z^{\ell} \bm \Pi_k.
\end{equation}
请注意，在上述公式中，矩阵 $\bm E^{\ell}$ 仅依赖于 $\Z^{\ell}$，其旨在{\em 扩展}所有特征以增加整体编码率；矩阵 $\bm C^{\ell}_{k}$ 依赖于第 $k$ 类的特征，旨在{\em 压缩}它们以降低每个类的编码率。
那么完整的梯度 $\frac{\partial \Delta
R_{\epsilon}}{\partial \bm Z}(\Z^\ell) \in \Re^{d\times N}$ 的形式为：
\begin{equation}
    \frac{\partial \Delta R_{\epsilon}}{\partial \bm Z}(\Z^\ell)  =
    \underbrace{\bm E^{\ell}}_{\text{Expansion}} \Z^{\ell} \;-\;
    \sum_{k=1}^K \gamma_{k} \underbrace{\bm
    C_{k}^{\ell}}_{\text{Compression}}  \Z^{\ell} \bm{\Pi}_k.
    \label{eqn:DR-gradient}
\end{equation}

\begin{remark}[将 $\bm E^\ell$ 和 $\bm C_j^\ell$ 解释为线性算子]\label{rem:regression-interpretation}
    对于任意 $\bm z^\ell \in \mathbb{R}^d$，
    \begin{gather}
        \bm E^\ell \bm z^\ell = \alpha(\bm z^\ell - \bm Z^\ell \bm
        q^\ell_\star), \qquad
        \mbox{where} \qquad \bm q^\ell_\star \doteq \argmin_{\bm
        q^\ell} \big\{\alpha \|\bm z^\ell - \bm Z^\ell \bm
        q^\ell\|_2^2 + \|\bm q^\ell\|_2^2\big\}.
    \end{gather}
    请注意，$\bm q^\ell_\star$ 恰好是所有相关数据点 $\bm Z^\ell$ 进行岭回归的解。
    因此，$\bm E^\ell$（对于 $\bm C^\ell_k$ 类似）近似于（即当 $N$ 足够大时）向由 $\bm Z^\ell$ 的列张成的子空间的正交补上的投影。解释矩阵 $\bm E^\ell$ 的另一种方法是通过协方差矩阵 $\Z^\ell (\Z^\ell)^\top$ 的特征值分解。假设 $\Z^\ell (\Z^\ell)^\top
    \doteq \bm U^\ell \bm \Lambda^\ell (\bm U^\ell)^\top$，其中 $\bm
    \Lambda^\ell \doteq \diag\left(\lambda^\ell_1, \ldots,
    \lambda^\ell_d \right)$，我们有
    \begin{equation}
        \bm E^\ell = \alpha \bm U^\ell\,
        \diag\left(\frac{1}{1+\alpha\lambda^\ell_1}, \ldots,
        \frac{1}{1+\alpha\lambda^\ell_d}\right) \left(\bm U^\ell\right)^\top.
    \end{equation}
    因此，矩阵 $\bm E^\ell$ 作用于向量 $\bm
    z^\ell$ 的方式是进行拉伸，使得大方差的方向被缩小，而方差消失的方向被保留。这些正是我们在 \eqref{eqn:expand-directions} 中移动特征的方向，以便整体体积扩大并且编码率增加，因此符号为正。产生相反效果的是，与 \eqref{eqn:compress-directions} 相关的方向是每个类的特征偏离其应属子空间的“残差”。这些正是特征需要被压缩回其各自子空间的方向，因此符号为负（见\Cref{fig:regression-interpretation}）。

    \begin{figure}[t]
        \centering
        \includegraphics[width=0.65\linewidth]{\toplevelprefix/chapters/chapter5/figs/expand_compress.png}
        \caption{\small $\bm C^\ell_k$ 和 $\bm
            E^\ell$ 的解释：$\bm C^\ell_k$ 通过将特征收缩到低维子空间来压缩每个类；$\bm E^\ell$ 通过对比和排斥不同类之间的特征来扩展所有特征。}
        \label{fig:regression-interpretation}
        \vspace{-0.1in}
    \end{figure}

    本质上，用于编码率衰减的梯度上升中的线性操作 $\bm E^\ell$ 和 $\bm
    C_k^\ell$ 是由执行“自回归”的训练数据决定的。最近在过参数化设定下对岭回归的重新认识 \cite{yang2020rethinking,Wu2020OnTO} 表明，使用看似冗余采样的数据（来自每个子空间）作为回归量不会导致过拟合。
\end{remark}

\paragraph{梯度引导的特征映射增量。} 请注意，在上述内容中，梯度上升将所有特征 $\Z^{\ell} =
[\z^{\ell}_{1}, \dots, \z^{\ell}_{N}]$ 视为自由变量。增量 $\Z^{\ell+1} - \Z^{\ell} = \eta \frac{\partial \Delta
R_{\epsilon}}{\partial \bm Z}(\Z^\ell)$ 尚未给出整个特征域 $\z^\ell \in \Re^d$ 上的变换。根据方程 \eqref{eqn:DR-gradient}，无法在成员归属未知的点处计算梯度，如\Cref{fig:gradient-flow}所示。因此，为了显式地找到最优的 $f(\x,\bm  \theta)$，我们可以考虑在第 $\ell$ 层特征 $\z^\ell$ 上构建一个小增量变换 $g(\cdot, \theta^{\ell})$，以模拟上述（投影）梯度方案：
\begin{equation}
    \z^{\ell + 1}   \; \propto \; \z^{\ell} + \eta\cdot  g(\z^{\ell},
    \theta^{\ell}) \quad \mbox{subject to} \quad \z^{\ell +1}
    \in \mathbb{S}^{d-1}
    \label{eqn:gradient-descent-transform}
\end{equation}
使得 $\big[g(\z_1^{\ell}, \bm \theta^{\ell}), \ldots,
g(\z_N^{\ell}, \bm \theta^{\ell}) \big] \approx \frac{\partial \Delta
R_{\epsilon}}{\partial \bm Z}(\Z^\ell).$ 也就是说，我们需要用定义在整个特征空间 $\z^\ell \in \Re^d$ 上的连续映射 $g(\z,\bm
\theta)$ 来近似局部形变所有（训练）特征 $\{\z_i^\ell\}_{i=1}^N$ 的梯度流 $\frac{\partial \Delta
R_{\epsilon}}{\partial \bm Z}$。请注意，人们可以将增量 \eqref{eqn:gradient-descent-transform} 解释为连续微分方程的离散化版本：
\begin{equation}
    \dot{\z} = g(\z, \theta).
\end{equation}
因此，如此构建的（深度）网络可以被解释为某种神经常微分方程 (Neural ODE) \cite{chen2018neural}。然而，与神经 ODE 中将流 $g$ 选择为某些通用结构不同，这里的 $g(\z, \bm\theta)$ 是为了模拟特征集上编码率衰减的梯度流（如\Cref{fig:gradient-flow}所示）：
\begin{equation*}
    \dot{\Z} = \frac{\partial \Delta R_{\epsilon}}{\partial \bm Z},
\end{equation*}
并且其结构完全由该目标推导并完全确定，没有任何其他先验或启发式方法。

通过检查梯度 \eqref{eqn:DR-gradient} 的结构，它表明增量变换 $g(\z^\ell, \bm \theta^\ell)$ 的一个自然候选形式为：
\begin{equation}
    g(\z^\ell, \bm \theta^\ell) \; \doteq \; \bm E^\ell \z^\ell -
    \sum_{k=1}^K \gamma_{k} \pi_k(\z^\ell)\bm C_{k}^{\ell}  \z^{\ell}
    \in \Re^d,
    \label{eqn:DR-gradient-transform}
\end{equation}
其中 $\pi_k(\z^\ell) \in [0,1]$ 表示 $\z^{\ell}$ 属于第 $k$ 类的概率。增量映射参数 $\bm \theta^\ell$ 依赖于：首先，由 $\bm E^{\ell}$ 和 $\{ \bm C^{\ell}_{k}\}_{k=1}^{K}$ 表示的一组线性映射，它们仅依赖于训练特征 $\Z^\ell$ 的统计量；其次，任意特征 $\z^\ell$ 的成员归属 $\{ \pi_k(\z^\ell)\}_{k=1}^K$。请注意，在成员归属 $\bm \Pi_k$ 已知的训练样本 $\Z^\ell$ 上，如此定义的 $g(\z^\ell, \bm
\theta)$ 恰好给出了梯度 $\frac{\partial \Delta
R_{\epsilon}}{\partial \bm Z}(\Z^\ell)$ 的值。

由于我们只有训练样本的成员归属，因此 \eqref{eqn:DR-gradient-transform} 中定义的函数 $g(\cdot)$ 只能在训练集上进行评估。为了将 $g(\cdot)$ 外推到整个特征空间，我们需要估计其第二项中的 $\pi_k(\z^\ell)$。在传统的深度学习中，该映射通常被建模为深度网络，并从训练数据中学习，例如通过{\em 反向传播}。然而，我们这里的目标并不是要学习一个精确的分类器 $\pi_{k}(\z^\ell)$。相反，我们只需要对类别信息进行足够好的估计，以便 $g(\cdot)$ 能够很好地近似梯度 $\frac{\partial \Delta
R_{\epsilon}}{\partial \bm Z}$。

根据\Cref{rem:regression-interpretation}给出的线性映射 $\bm E^\ell$ 和 $\bm C_k^\ell$ 的几何解释，项 $\bm p_{k}^{\ell} \doteq \bm C^{\ell}_k \z^{\ell}$ 可以被视为（近似地）$\z^{\ell}$ 在每个类 $j$ 的正交补上的投影。因此，如果 $\z^\ell$ 属于类 $j$，则 $\|\bm p_{j}^{\ell}\|_2$ 较小，否则较大。这启发我们基于以下 softmax 函数来估计其成员归属：
\begin{eqnarray}
    \widehat{\vpi}(\z^\ell) &\doteq&
    \softmax\rp{-\lambda\mat{\norm{\vC^{\ell}_{1}\vz^{\ell}}_{2}
    \\ \vdots \\ \norm{\vC^{\ell}_{K}\vz^{\ell}}_{2}}}\nonumber \\
    &=& \frac{1}{\sum_{k = 1}^{K}\exp(-\lambda
    \norm{\vC^{\ell}_{k}\vz^{\ell}}_{2})}\mat{\exp(-\lambda\norm{\vC^{\ell}_{1}\vz^{\ell}}_{2})
    \\ \vdots \\ \exp(-\lambda\norm{\vC^{\ell}_{K}\vz^{\ell}}_{2})}
    \in [0,1]^{K}.
\end{eqnarray}
因此，\eqref{eqn:DR-gradient-transform} 的第二项可以通过这个估计的成员归属来近似：
\begin{align}
    \sum_{k=1}^K \gamma_{k} \pi_k(\z^\ell)\bm C_{k}^{\ell}  \z^{\ell}
    \; \approx \;  \sum_{k=1}^K \gamma_{k} \widehat{\pi}_k(\z^\ell)
    \bm{C}^{\ell}_k  \z^{\ell}
    \; \doteq \; \bm \sigma\Big([\bm{C}^{\ell}_{1} \z^{\ell}, \dots,
    \bm{C}^{\ell}_{K} \z^{\ell}]\Big),
    \label{eqn:soft-residual}
\end{align}
这被表示为特征 $\z^\ell$ 通过 $K$ 组滤波器 $[\bm{C}^{\ell}_{1}, \dots, \bm{C}^{\ell}_{K}]$ 的输出上的非线性算子 $\bm \sigma(\cdot)$。请注意，非线性的产生是由于基于这些滤波器的特征响应对类别成员归属进行了“软”分配。

总而言之，结合 \eqref{eqn:gradient-descent-transform}、\eqref{eqn:DR-gradient-transform} 和 \eqref{eqn:soft-residual}，从 $\z^{\ell}$ 到 $\z^{\ell+1}$ 的增量特征变换现在变为
\begin{equation}\label{eqn:layer-approximate}
    \begin{aligned}
        \z^{\ell+1}  &\propto \; \z^\ell +  \eta \cdot  \bm E^{\ell}
        \z^{\ell} - \eta\cdot  \bm \sigma\big([\bm{C}^{\ell}_{1}
        \z^{\ell}, \dots, \bm{C}^{\ell}_{K} \z^{\ell}]\big)  \\
        &= \; \z^\ell +  \eta \cdot g(\z^\ell, \bm \theta^\ell)
        \qquad \mbox{s.t.} \quad \z^{\ell +1} \in \mathbb{S}^{d-1},
    \end{aligned}
\end{equation}
其中非线性函数 $\bm \sigma(\cdot)$ 如上定义，$\bm \theta^\ell$ 收集了所有逐层参数。即 $\bm \theta^\ell =\left\{\bm E^\ell, \bm{C}^{\ell}_{1}, \dots,
\bm{C}^{\ell}_{K}, \gamma_{k}, \lambda\right\}$。请注意，每一层的特征总是通过投影到单位球面 $\mathbb S^{d-1}$ 上进行“归一化”，记为 $\mathcal P_{\mathbb S^{d-1}}$。\eqref{eqn:layer-approximate} 中增量的形式可以通过\Cref{fig:arch-ReduNet}(a)中的图表来说明。

\begin{figure}[t]
    \begin{subfigure}[t]{0.4\textwidth}
        \centering
        \includegraphics[width=\textwidth]{\toplevelprefix/chapters/chapter5/figs/redunet_layer.png}
        \caption{\textbf{ReduNet}}
    \end{subfigure}
    \hfill
    \begin{subfigure}[t]{0.58\textwidth}
        \centering
        \includegraphics[width=\textwidth]{\toplevelprefix/chapters/chapter5/figs/resnet_resnext.pdf}
        \caption{\textbf{ResNet} and \textbf{ResNeXt}.}
    \end{subfigure}
    \caption{\small ReduNet 的网络架构及与其他网络的比较。\textbf{(a)}：由优化编码率衰减的梯度上升的一次迭代推导出的 \textbf{ReduNet} 层结构。\textbf{(b)}（左）：ResNet~\cite{he2016deep} 的一层；以及 \textbf{(b)}（右）：ResNeXt~\cite{ResNEXT} 的一层。正如我们将在\Cref{sec:shift-invariant}中看到的，当施加平移不变性时，ReduNet 的线性算子 $\bm E^\ell$ 和 $\bm{C}_k^\ell$ 自然会变成（多通道）卷积。}
    \label{fig:arch-ReduNet}
\end{figure}

\begin{algorithm}[t]
    \caption{ReduNet 的训练算法}
    \label{alg:redunet_training}
    \begin{algorithmic}[1]
        \Require{$\vX = [\vx_1, \ldots, \vx_N]\in \Re^{D \times N}$，$\bm{\Pi} = \{\bm \Pi_k\}_{k=1}^K$，$\epsilon > 0$，$\lambda$，以及学习率 $\eta$。}
        \Ensure{学习到的参数 \(\{\vE^{\ell}\}_{\ell =
                1}^{L}, \{\{\vC^{\ell}_{k}\}_{k = 1}^{K}\}_{\ell = 1}^{L},
        \{\gamma_{k}\}_{k = 1}^{k}\)}

        \Procedure{ReduNetTraining}{$\vX, \vPi, \epsilon, \lambda, \eta$}
        \Statex{\quad\ \texttt{\# 定义常量}}

        \State{\(\alpha \gets d/(N\epsilon^{2})\)}
        \For{\(k \in \{1, \dots, K\}\)}
        \State{\(\alpha_{k} \gets D/(\tr(\vPi_{k})\epsilon^{2})\)}
        \State{\(\gamma_{k} \gets \tr(\vPi_{k})/D\)}
        \EndFor
        \Statex{}

        \Statex{\quad\ \texttt{\# ReduNet 逐层迭代}}
        \State{\(\vZ^{1} = \mat{\vz_{1}^1, \dots, \vz_{N}^1} \gets
        \vX\)} \Comment{初始化 ReduNet 逐层迭代}
        \For{\(\ell \in \{1, \dots, L\}\)}
        \Statex{\quad\ \texttt{\# 步骤 1：计算网络参数}
        \(\vE^{\ell}, \{\vC^{\ell}_{k}\}_{k = 1}^{K}\)}
        \State{\(\vE^{\ell} \gets \alpha\left(\vI + \alpha
        \vZ^{\ell}(\vZ^{\ell})^{\top}\right)^{-1} \in \R^{d \times d}\)}
        \For{\(k \in \{1, \dots, K\}\)}
        \State{\(\vC^{\ell}_{k} \gets \alpha_{k}\left(\vI +
                \alpha_{k}\vZ^{\ell}\vPi_{k}(\vZ^{\ell})^{\top}\right)^{-1}
        \in \R^{d \times d}\)}
        \EndFor
        \Statex{}
        \Statex{\qquad\quad\texttt{\# 步骤 2：更新特征 \(\vZ^{\ell}\)}}
        \For{\(i \in \{1, \dots, N\}\)}
        \State{\(\hat{\vpi}(\vz^{\ell}_i) \gets \displaystyle
                \softmax(-\lambda[\norm{\vC^{\ell}_{1}\vz^{\ell}_i}_{2},
                \dots, \norm{\vC^{\ell}_{K}\vz^{\ell}_i}_{2}]) \in [0,
        1]^{K}\)} \Comment{计算软分配
        \(\hat{\vpi}(\vz^{\ell}_i)\)}
        \State{\(\displaystyle\vz^{\ell + 1}_i \gets \cP_{\bS^{d -
                1}}\rp{\vz^{\ell}_i + \eta \bp{\vE^{\ell}\vz^{\ell}_i -
                        \sum_{k = 1}^{K}\gamma_{k}\hat{\pi}_{k}(\vz^{\ell}_i)
        \vC^{\ell}_{k}\vz^{\ell}_i}} \in \R^{d}\)} \Comment{从 \(\vz^{\ell}_i\) 更新特征 \(\vz^{\ell + 1}_i\)}
        \EndFor
        \EndFor
        \State{\Return{\(\{\vE^{\ell}\}_{\ell = 1}^{L},
                    \{\{\vC^{\ell}_{k}\}_{k = 1}^{K}\}_{\ell = 1}^{L},
        \{\gamma_{k}\}_{k = 1}^{K}\)}} \Comment{返回所有网络参数。}
        \EndProcedure
    \end{algorithmic}
\end{algorithm}

\paragraph{用于优化编码率衰减的深度网络。} 请注意，增量的构建是为了模拟编码率衰减 $\Delta R_\epsilon$ 的梯度上升。因此，通过上述过程迭代地转换特征，我们期望编码率衰减会增加，正如我们将在实验部分看到的那样。这个迭代过程一旦收敛（例如在 $L$ 次迭代之后），就会在输入 $\x = \z^0$ 上给出所需的特征映射 $f(\x, \bm \theta)$，其形式恰好是一个{\em 深度网络}，其中每一层具有\Cref{fig:arch-ReduNet}左侧所示的结构：
\begin{equation}\label{eqn:ReduNet}
    \begin{aligned}
        f(\x, \bm \theta)\; =&  \;\;f^L \circ f^{L-1} \circ  \cdots
        \circ f^1 \circ
        f^0(\z^0),  \\
        f^\ell(\z^\ell, \bm \theta^\ell) \; \doteq & \;\; \z^{\ell+1}
        = \mathcal{P}_{\mathbb{S}^{n-1}}[\z^{\ell} + \eta\cdot
        g(\z^{\ell}, \bm \theta^{\ell})], \\
        g(\z^{\ell}, \bm \theta^{\ell}) \; =&\;\; \bm E^{\ell}
        \z^{\ell} -  \bm \sigma\big([\bm{C}^{\ell}_{1} \z^{\ell},
        \dots, \bm{C}^{\ell}_{K} \z^{\ell}]\big).
    \end{aligned}
\end{equation}
由于这个深度网络源于最大化编码率衰减 (rate \textbf{redu}ction)，我们将其称为 \textbf{ReduNet}。通过将 ReduNet 的架构与\Cref{fig:arch-ReduNet}中所示的流行的凭经验设计的网络 \textbf{ResNet} 和 \textbf{ResNeXt} 的架构进行比较，其相似性有些不可思议。从概念上讲，ReduNet 也可以用来证明流行的混合专家 (\textbf{MoE}) 架构 \cite{MoE} 的合理性，因为每个并行通道 $\bm{C}^{\ell}_k$ 都可以被视为针对每类对象训练的“专家”。

\begin{figure}[t]
    \centering
    \includegraphics[width=0.45\linewidth]{\toplevelprefix/chapters/chapter5/figs/MoE.png}
    \hspace{5mm}
    \includegraphics[width=0.45\linewidth]{\toplevelprefix/chapters/chapter5/figs/switched-transformer.png}
    \caption{左：混合专家 (MoE) 深度网络 \cite{MoE}。右：促进稀疏性的 Switch Transformer \cite{Fedus-2022}，用于实现具有 1.7 万亿参数的 MoE。}
    \label{fig:enter-label}
\end{figure}

我们分别在\Cref{alg:redunet_training}和\Cref{alg:redunet_evaluation}中总结了 ReduNet 的训练和评估。请注意，网络的所有参数都是以{\em 前向传播}的方式逐层显式构建的。该构建过程不需要任何反向传播！如此学习到的特征可以直接用于分类，例如通过最近子空间分类器。

\begin{algorithm}[!htbp]
    \caption{ReduNet 的评估算法}
    \label{alg:redunet_evaluation}
    \begin{algorithmic}[1]
        \Require{输入 \(\vx \in \R^{D}\)，网络参数 \(\{\vE^{\ell}\}_{\ell = 1}^{L}, \{\{\vC^{\ell}_{k}\}_{k =
            1}^{K}\}_{\ell = 1}^{L}, \{\gamma_{k}\}_{k = 1}^{K}\)，学习率 \(\lambda\)}
        \Ensure{特征 \(\vz^{L + 1}\)}

        \Procedure{ReduNetEvaluation}{$\vx$}
        \State{\(\vz^{1} \gets \vx \in \R^{D}\)} \Comment{初始化 ReduNet 逐层迭代}
        \For{\(\ell \in \{1, \dots, L\}\)}
        \State{\(\hat{\vpi}(\vz^{\ell}) \gets
                \softmax(-\lambda\mat{\norm{\vC^{\ell}_{1}\vz^{\ell}}_{2},
                \dots, \norm{\vC^{\ell}_{K}\vz^{\ell}}_{2}}) \in [0,
        1]^{K}\)} \Comment{计算软分配 \(\hat{\vpi}(\vz^{\ell})\)}
        \State{\(\vz^{\ell + 1} \gets \cP_{\bS^{d - 1}}\rp{\vz^{\ell}
                    + \eta \bp{\vE^{\ell}\vz^{\ell} - \sum_{k =
                        1}^{K}\gamma_{k}\hat{\pi}_{k}(\vz^{\ell})
        \vC^{\ell}_{k}\vz^{\ell}}} \in \R^{d}\)} \Comment{使用 \(\vz^{\ell}\) 更新特征 \(\vz^{\ell + 1}\)}
        \EndFor
        \State{\Return{\(\vz^{L + 1}\)}} \Comment{返回输出特征}
        \EndProcedure
    \end{algorithmic}
\end{algorithm}

\begin{example}
    为了提供关于 ReduNet 如何转换特征的一些直观理解，我们提供了一个混合 3D 高斯分布的简单示例，并在\Cref{fig:redu-3d-gaussian-diagram}中可视化了特征是如何被转换的。考虑 $\R^{3}$ 中投影到 $\mathbb{S}^2$ 上的三个高斯分布的混合。我们首先为 3 个类生成数据点：对于 $k=1,2,3$，$\bm{X}_{k}=[\bm{x}_{k,1},
    \ldots, \bm{x}_{k,m}] \in \R^{3\times m}$，$\bm{x}_{k,i} \sim
    \mathcal{N}(\bm{\mu}_{k}, \sigma_{k}^{2} \I)$，且 ${\pi}(\x_{k,i}) = k$。我们设置 $m=500, \sigma_{1}=\sigma_{2}=\sigma_{3}=0.1$，并且 $\bm{\mu}_{1}, \bm{\mu}_{2}, \bm{\mu}_{3} \in \mathbb{S}^2$。然后我们将所有数据点投影到 $\mathbb{S}^{2}$ 上，即 $\bm{x}_{k,i}/\|\bm{x}_{k,i}\|_{2}$。为了构建网络（计算第 $\ell$ 层的 $\bm{E}^{\ell},
    \bm{C}^{\ell}_{k}$），我们设置迭代次数/层数 $L=2,000$，步长 $\eta=0.5$，精度 $\epsilon=0.1$。我们这样做只是为了证明我们的框架即使在数千层的情况下也能产生稳定的深度网络。在实践中，可能不需要数千层，当添加新层带来的收益递减时，就可以停止。对于这个例子，几百层就足够了。因此，清晰的优化目标为所需网络的深度提供了一个自然的准则。

    如\Cref{fig:redu-3d-gaussian-diagram}所示，我们可以观察到，在映射 $f(\cdot, \theta)$ 之后，来自同一类的样本被高度压缩并收敛到单个簇，并且两个不同簇之间的角度约为 $\pi/2$，这与 $\mathbb{S}^2$ 中 MCR$^2$ 损失的最优解 $\Z^{\star}$ 非常吻合。不同层上特征的 MCR$^2$ 损失可以在\Cref{fig:redu-3d-gaussian-diagram}(\textbf{c})中找到。经验上，我们发现构建的 ReduNet 能够最大化 MCR$^2$ 损失并稳定收敛，同一类的样本收敛到一个簇，不同簇彼此正交。此外，当从相同分布中采样新数据点时，我们发现来自同一类的新样本始终收敛到与训练样本相同的簇中心。
    \begin{figure}[t]
        \begin{subfigure}[t]{0.32\textwidth}
            \centering
            \includegraphics[width=\textwidth]{\toplevelprefix/chapters/chapter5/figs/scatter3d-X_train.pdf}\vspace{-0.1in}
            \caption{$\bm{X}_{\text{train}}$}
        \end{subfigure}
        \hfill
        \begin{subfigure}[t]{0.32\textwidth}
            \centering
            \includegraphics[width=\textwidth]{\toplevelprefix/chapters/chapter5/figs/scatter3d-Z_train.pdf}\vspace{-0.1in}
            \caption{$\bm{Z}_{\text{train}}$}
        \end{subfigure}
        \hfill
        \begin{subfigure}[t]{0.32\textwidth}
            \centering
            \includegraphics[width=\textwidth]{\toplevelprefix/chapters/chapter5/figs/scatter3d-loss-traintest.pdf}\vspace{-0.1in}
            \caption{Loss (train/val)}
        \end{subfigure}
        \vspace{-0.1in}
        \caption{\small 3D 高斯混合的原始样本和学习到的表示。我们通过散点图在 (\textbf{a}) 中可视化了数据点 $\bm{X}$（映射 $f(\cdot, \theta)$ 之前），在 (\textbf{b}) 中可视化了学习到的特征 $\bm{Z}$（映射 $f(\cdot, \theta)$ 之后）。在每个散点图中，每种颜色代表一类样本。在 (\textbf{c}) 中，我们还展示了目标函数值随时间变化的曲线图。}
        \label{fig:redu-3d-gaussian-diagram}
    \end{figure}

\end{example}

\subsection{源自不变编码率衰减的卷积网络}\label{sec:shift-invariant}

在上一节中，我们使用针对编码率衰减目标的展开优化推导了深度网络 ReduNet 的逐层架构。具体而言，\eqref{eq:mcr-formulation} 中的压缩项 $R^c_\epsilon(\Z \mid\bm \Pi)$ 旨在压缩来自同一类的表示。然而，该公式没有考虑输入数据可能的域变换或形变。例如，将对象稍微向右平移不会改变图像的语义标签。在本节中，我们将演示如何通过最大化对某些域形变（如图像旋转和平移）具有不变性的编码率衰减目标来推导卷积层。

对于许多聚类或分类任务（例如图像中的目标检测），如果两个样本相差某些类别的域形变或数据增强 $\cT = \{\tau \}$，我们认为它们是{\em 等价的}。因此，我们只对对此类形变具有{\em 不变性}的低维结构感兴趣（即 $\x \in \mathcal{M}$ 当且仅当对于所有 $\tau \in \cT$ 都有 $\tau(\x) \in \mathcal{M}$），已知这些结构具有复杂的几何和拓扑结构，在实践中很难精确学习，即使使用严格设计的 CNN 也是如此 \cite{Cohen-ICML-2016}。在这个框架中，这可以以一种非常自然的方式表述：所有等变实例都将被嵌入到同一个子空间中，从而使得该子空间本身对所考虑的变换具有不变性。

在许多应用中，例如序列数据或图像数据，数据的语义含义（标签）对某些变换 $\mathfrak{g} \in \mathbb{G}$（对于某个群 $\mathbb{G}$）是{\em 不变的} \cite{CohenW16,deep-sets-NIPS2017}。例如，音频信号的含义对时间平移是不变的；图像中对象的身份对图像平面中的平移是不变的。因此，我们希望特征映射 $f(\x,\bm \theta)$ 严格地对此类变换保持不变：
\begin{equation}
    \mbox{\em 群不变性：}\;   f(\x\circ \mathfrak{g}, \bm
    \theta) \sim f(\x,\bm \theta),\ \forall \mathfrak{g} \in \mathbb{G},
\end{equation}
其中“$\sim$”表示两个特征属于同一个等价类。尽管为了确保不变性或等变性，卷积算子已成为深度网络中的常见做法 \cite{CohenW16}，但在实践中，从头开始训练一个（凭经验设计的）卷积网络，使其能够{\em 保证}即使对平移和旋转等简单变换也具有不变性，仍然具有挑战性 \cite{azulay2018deep,engstrom2017rotation}。另一种方法是仔细设计每一层的卷积滤波器，以确保对各种信号的平移不变性，例如像 ScatteringNet \cite{scattering-net} 及其后续工作 \cite{Wiatowski-2018} 中那样使用小波。然而，为了确保对通用信号的不变性，所需的卷积数量通常随网络深度呈指数增长。这就是为什么这种类型的网络不能构建得很深，通常只有几层的原因。

现在，我们展示 MCR$^2$ 原则以一种自然而精确的方式与不变性兼容：我们只需要将所有变换版本 $\{\x\circ \mathfrak{g} \mid \mathfrak{g} \in \mathbb G\}$ 分配到与数据 $\x$ 相同的类中，并将它们的特征 $\z$ 全部映射到同一个子空间 $\mathcal S$。因此，所有群等变信息仅编码在子空间内部，并且在所得子空间集合上定义的任何分类器都将自动对此类群变换具有不变性。有关 1D 旋转和 2D 平移示例的说明，请参见\Cref{fig:ortho-invariance-diagram}。接下来，我们将严格证明，当群 $\mathbb G$ 是循环 1D 平移时，所得的深度网络自然会变成一个{\em 多通道卷积网络}。因为如此构建的网络只需要确保给定数据 $\X$ 或其特征 $\Z$ 的不变性，所以与 ScatteringNet 相反，所需的卷积数量在非常深的网络中实际上保持不变。

\begin{figure}[t]
    \begin{subfigure}[t]{0.4\textwidth}
        \centering
        \includegraphics[width=\textwidth]{\toplevelprefix/chapters/chapter5/figs/ortho_diagram_1d.pdf}
    \end{subfigure}
    \hfill
    \begin{subfigure}[t]{0.4\textwidth}
        \centering
        \includegraphics[width=\textwidth]{\toplevelprefix/chapters/chapter5/figs/ortho_diagram_2d.pdf}
    \end{subfigure}
    \caption{对图像旋转（左）或平移（右）具有等变性/不变性的所求表示的图示：每个类的所有变换图像都被映射到同一个与其他子空间不相干的子空间中。嵌入在每个子空间中的特征对变换群是等变的，而每个子空间对此类变换是不变的。}\label{fig:ortho-invariance-diagram}
\end{figure}

\paragraph{1D 序列数据与平移不变性} 为了对在平移下不变的一维数据 $\bm x = [x(0), x(1), \ldots, x(D-1)] \in \Re^D$ 进行分类，我们将 $\mathbb{G}$ 视为所有循环平移的群。每个观测值 $\bm x_i$ 生成一个平移副本族 $\{
\x_i \circ \mathfrak{g} \, | \, \mathfrak{g} \in \mathbb G \}$，它们是循环矩阵 $\circm(\bm x_i) \in \Re^{D \times D}$ 的列，由下式给出：
\begin{equation}
    \circm(\x) \,\doteq\, \left[
        \begin{array}{ccccc} x(0) & x(D-1) & \dots & x(2) & x(1)
            \\ x(1) & x(0) & x(D-1) & \cdots & x(2) \\ \vdots & x(1)
            & x(0) &\ddots & \vdots \\ x(D-2) &  \vdots & \ddots &
            \ddots & x(D-1) \\ x(D-1) & x(D-2) & \dots & x(1) & x(0)
    \end{array} \right]  \in \Re^{D \times D}.
\end{equation}
关于循环矩阵的性质，我们建议读者参考 \cite{Kra2012OnCM}。为简单起见，令 $\bm Z^1 \doteq [ \z_{1}^1, \dots,
\z_{N}^1 ] = \X \in \Re^{d \times N}$。\footnote{同样，为了简化讨论，我们暂时假设初始特征 $\Z^1$ 就是 $\X$ 本身，因此具有相同的维度 $d$，即 $D=d$。但这并非必须如此，因为我们很快就会看到我们需要将 $\X$ 提升到更高的维度。}那么，如果我们从它们的循环族 $\circm(\bm Z^1) = \left[
\circm(\z_{1}^1), \dots, \circm(\z_{N}^1) \right] \in \Re^{d \times
(dN)}$ 构建 ReduNet 会发生什么？也就是说，我们希望通过 ReduNet 将所有这些压缩并映射到同一个子空间中。

请注意，现在数据协方差矩阵：
\begin{eqnarray}
    \circm(\bm Z^1) \circm(\bm Z^1)^\top
    &=& \left[ \circm(\z_{1}^1), \dots, \circm(\z_{N}^1) \right]
    \left[ \circm(\z_{1}^1), \dots, \circm (\z_{N}^1) \right]^\top \nonumber \\
    &=& \sum_{i =1}^N \circm(\z_{i}^1) \circm(\z_{i}^1)^\top \;\in
    \Re^{d\times d}
\end{eqnarray}
与该样本族相关的矩阵{\em 自动}是一个（对称）循环矩阵。此外，由于循环性质在求和、求逆和乘积运算下保持不变，因此矩阵 $\bm E^1$ 和 $\bm C^1_k$ 也自动是循环矩阵，它们对特征向量 $\bm z \in \Re^d$ 的应用可以使用循环卷积“$\circledast$”来实现。具体来说，我们有以下命题。

\begin{proposition}[$\bm E^1$ 和 $\bm C^1_k$ 的卷积结构]
    矩阵
    \begin{equation}
        \bm E^1 = \alpha\big(\bm I + \alpha \circm(\bm Z^1)
        \circm(\bm Z^1)^\top \big)^{-1}
    \end{equation}
    是一个循环矩阵，表示一个循环卷积：
    $$\bm E^1 \z = \bm e_1 \circledast \z,$$
    其中 $\bm e_1 \in \Re^d$ 是 $\bm E^1$ 的第一列向量，“$\circledast$”是定义如下的循环卷积：
    \begin{equation*}
        (\bm e_1 \circledast \bm z)_{i} \doteq \sum_{j=0}^{d-1}
        e_1(j) x(i+ d-j \,\, \textsf{mod} \,\,d).
    \end{equation*}
    类似地，与 $\bm Z^1$ 的任何子集相关的矩阵 $\bm C^1_k$ 也是循环卷积。
    \label{prop:circular-conv-1}
\end{proposition}

不仅 ReduNet 的第一层参数 $\bm E^1$ 和 $\bm{C}^1_k$ 变成了循环卷积，而且下一层特征也保持为循环矩阵。也就是说，应用于 $\z^1 \in \Re^d$ 的所有平移版本的 \eqref{eqn:layer-approximate} 中的增量特征变换，由下式给出：
\begin{equation}
    \circm(\z^1) + \eta \cdot \bm E^1 \circm(\z^1) - \eta \cdot \bm
    \sigma \Big([\bm{C}_1^1 \circm(\z^1), \ldots, \bm{C}^1_K
    \circm(\z^1)] \Big),
\end{equation}
是一个循环矩阵。这意味着不需要像第一层那样从第二层特征构建循环族。通过记作
\begin{equation}
    \z^{2} \propto \z^{1} +\eta \cdot g(\z^{1}, \bm \theta^{1}) =
    \z^1 + \eta \cdot \bm e_{1} \circledast \z^{1} -  \eta \cdot \bm
    \sigma\Big([\bm{c}_{1}^{1} \circledast \z^{1}, \dots,
    \bm{c}^{1}_{K} \circledast \z^{1}]\Big),
    \label{eqn:approximate-convolution}
\end{equation}
下一层的特征可以写为
$$\circm(\bm Z^2) = \big[ \circm( \bm z_{1}^1 + \eta g( \bm z_{1}^1,
    \bm \theta^1)), \dots, \circm( \bm z_{N}^1 + \eta g(\bm z_{N}^1, \bm
\theta^1)) \big].$$
通过归纳，我们可以看到所有基于这种 $\circm(\bm Z^\ell)$ 的矩阵 $\bm E^\ell$ 和 $\bm C^\ell_k$ 都是循环的，所有特征也是如此。凭借数据的性质，ReduNet 已经呈现出卷积网络的形式，{\em 而无需显式地选择这种结构！}

\paragraph{不变性与稀疏性之间的基本权衡。} 但存在一个问题：通常，向量 $\z$ 的所有循环排列的集合给出了一个满秩矩阵。也就是说，与每个样本（因此每个类）相关的 $d$ 个“增强”特征通常已经张成了整个空间 $\Re^d$。例如，delta 函数 $\delta(d)$ 的所有平移版本可以生成任何其他信号作为它们的（稠密）加权叠加。MCR$^2$ 目标 \eqref{eqn:maximal-rate-reduction} 将无法将类区分为不同的子空间。

一种自然的补救措施是通过将原始信号“提升”到更高维的空间来提高数据的可分离性，例如，通过获取它们对多个滤波器 $\bm k_1, \ldots, \bm k_C \in \Re^d$ 的响应：
\begin{equation}
    \bm z[c] = \bm k_c \circledast \bm x  =  \circm(\bm k_c) \bm x
    \in \Re^d, \quad c = 1, \ldots, C.
    \label{eqn:lift-1d}
\end{equation}
这些滤波器可以是预先设计的促进不变性的滤波器，\footnote{对于像音频这样的 1D 信号，可以考虑传统的短时傅里叶变换 (STFT)；对于 2D 图像，可以考虑像 ScatteringNet \cite{scattering-net} 中那样的 2D 小波。} 或者是从数据中自适应学习的，\footnote{对于学习到的滤波器，可以像 PCANet \cite{chan2015pcanet} 中那样将滤波器学习为样本的主成分，或者从卷积字典学习 \cite{li2019multichannel,qu2019nonconvex} 中学习。} 或者是像我们在实验中所做的那样随机选择的。此操作将每个原始信号 $\x \in \Re^d$ 提升为 $C$ 通道特征，记为 $\bar{\z}  \doteq
[\z[1], \ldots, \z[C]]^\top \in \Re^{C\times d}$。然后，我们可以在 $\bar{\z}$ 的向量表示上构建 ReduNet，记为 $\vec(\bar\z) \doteq [\z[1]^\top, \ldots, \z[C]^\top] \in \Re^{dC}$。相关的循环版本 $ \circm(\bar{\z})$ 及其数据协方差矩阵（记为 $\bar{\bm \Sigma}(\bar\z)$）对于其所有平移版本给出如下：
\begin{equation}
    \begin{aligned}
        \circm(\bar{\z}) &\doteq \left[
            \begin{matrix}
                \circm(\z[1])  \\ \vdots \\ \circm(\z[C])
        \end{matrix} \right] \in \Re^{dC\times d},
        \\
        \bar{\bm \Sigma}(\bar\z) &\doteq
        \left[
            \begin{matrix}
                \circm(\z[1]) \\ \vdots \\ \circm(\z[C])
        \end{matrix} \right]\left[
            \begin{matrix}\circm(\z[1])^\top,\ldots, \circm(\z[C])^\top
        \end{matrix} \right] \in \Re^{dC\times dC},
        \label{eqn:W-multichannel}
    \end{aligned}
\end{equation}
其中 $c \in [C]$ 的 $\circm(\z[c]) \in \Re^{d\times d}$ 是特征 $\bar \z$ 的第 $c$ 个通道的循环版本。那么 $\circm(\bar\z)$ 的列最多只能张成 $\Re^{dC}$ 中的一个 $d$ 维真子空间。然而，这种简单的提升操作（如果是线性的）还不足以使类可分离——与其他类相关的特征将张成{\em 相同的} $d$ 维子空间。这反映了不变性与线性（子空间）建模之间的基本冲突：{\em 人们不能指望任意平移和叠加的信号属于同一个类。}

\begin{figure}[t]
    \centerline{
    \includegraphics[width=0.9\textwidth]{\toplevelprefix/chapters/chapter5/figs/sparse-representation.pdf}}
    \caption{每个输入信号 $\bm x$（这里是一张图像）可以表示为与字典 $\bm D$ 中的多个核 $\bm d_c$ 的稀疏卷积的叠加。}
    \label{fig:multi-channel-sparse-representation}
\end{figure}
解决这种冲突的一种方法是利用每个类内部的额外结构，即{\em 稀疏性}的形式：每个类内的信号不是作为某些基础原子（或基序）的任意线性叠加生成的，而仅仅是它们及其平移版本的{\em 稀疏}组合，如\Cref{fig:multi-channel-sparse-representation}所示。更准确地说，令 $\bm D_k = [\bm d_{k,1}, \ldots, \bm d_{k,c}]$ 表示一个包含与类 $k$ 相关的原子集合的矩阵，也称为字典，那么该类中的每个信号 $\x$ 稀疏地生成为：
\begin{equation}
    \x = \bm d_{k,1} \circledast z_1 + \ldots + \bm d_{k,c}
    \circledast z_c = \circm(\bm{D}_k)\z,
\end{equation}
对于某个稀疏向量 $\z$。不同类中的信号则
由不同的字典生成，这些字典的原子（或基元）彼此之间是不相干的。由于这种不相干性，某一类的信号不太可能被任何其他类的原子稀疏地表示。因此，所有信号都可以表示为
\begin{equation}
    \x = \big[\circm(\bm{D}_1), \circm(\bm{D}_2), \ldots,
    \circm(\bm{D}_K)\big] \bar \z,
\end{equation}
其中 $\bar \z$ 是稀疏的。\footnote{请注意，类似的稀疏表示模型早被提出并用于人脸识别等应用的分类任务中，展现出了极好的效果 \cite{Wright:2009,wagner2012toward}。最近，\cite{papyan2017convolutional} 提出了卷积稀疏编码模型，作为解释深度卷积网络结构的框架。} 关于如何从样本数据中学习最紧凑和最优的稀疏化字典（例如 \cite{li2019multichannel,qu2019nonconvex}），并随后求解逆问题以计算相关的稀疏代码 $\z$ 或 $\bar \z$，已有大量文献。最近 \cite{qu2020nonconvex,Qu2020Geometric} 的研究甚至表明，在宽泛的条件下，卷积字典学习问题可以得到有效且高效的解决。

然而，对于分类等任务，我们不一定对精确的最优字典或每个单独信号的精确稀疏代码感兴趣。我们主要关心的是，每一类的稀疏代码集合作为一个整体，是否能与其他类的稀疏代码充分可分。在稀疏生成模型的假设下，如果卷积核 $\{\bm k_c\}_{c=1}^C$ 与上述稀疏化字典 $\bm D = [\bm D_1, \ldots, \bm D_K]$ 的“转置”或“逆”（也称为{\em 分析滤波器} \cite{Cosparse-Nam,Analysis-Filter}）匹配良好，那么某一类的信号将仅对这些滤波器的一小部分产生高响应，而对其他滤波器产生低响应（由于不相干性假设）。尽管如此，在实践中，通常足够数量（例如 $C$ 个）的随机滤波器 $\{\bm k_c\}_{c=1}^C$ 就足以确保提取的 $C$ 通道特征
\begin{equation}
    \big[\bm k_1 \circledast \x, \bm k_2 \circledast \x, \ldots, \bm
    k_C \circledast \x\big]^\top = \big[\circm(\bm k_1) \x, \ldots,
    \circm(\bm k_C) \x \big]^\top \in \Re^{C\times d}
\end{equation}
对于不同类别对不同滤波器具有不同的响应模式，从而使不同类别变得可分 \cite{chan2015pcanet}。

因此，在我们的框架中，通道数（或网络的宽度）在很大程度上真正扮演了{\em 统计资源}的角色，而层数（网络的深度）则扮演了{\em 计算资源}的角色。压缩感知理论精确地刻画了需要多少次测量才能保留数据的内在低维结构（包括可分性）\cite{Wright-Ma-2022}。 %

\begin{figure}[t]
    \centerline{
    \includegraphics[width=0.95\textwidth]{\toplevelprefix/chapters/chapter5/figs/sparse-lifting.pdf}}
    \caption{通过与多个核 $\bm k_c$ 进行卷积然后进行稀疏化，来估计输入信号 $\bm x$（此处为图像）的稀疏代码 $\bar{\bm z}$。}
    \label{fig:multi-channel-sparse-lifting}
\end{figure}
多通道响应 $\bar \z$ 应该是稀疏的。因此，为了近似稀疏代码 $\bar \z$，我们可以对上述滤波器输出采用逐元素的{\em 促进稀疏的非线性阈值化}，例如 $\bm \tau(\cdot)$，通过将较低的（例如绝对值低于 $\epsilon$）或负的响应设为零：
\begin{equation}
    \bar \z \doteq \bm \tau \left( \big[\circm(\bm k_1) \x, \ldots,
    \circm(\bm k_C) \x \big]^\top \right) \in \Re^{C \times d}.
    \label{eqn:sparse-lifting}
\end{equation}
\Cref{fig:multi-channel-sparse-lifting} 说明了基本思想。关于稀疏化阈值算子设计的更系统性研究，可以参考 \cite{Analysis-Filter}。然而，在这里我们对获得最佳稀疏代码并不那么感兴趣，只要这些代码足够可分即可。因此，非线性算子 $\bm \tau$ 可以简单地选择为软阈值化或 ReLU。这些假定为稀疏的特征 $\bar \z$ 可以被假设位于 $\mathbb{R}^{dC}$ 的一个低维（非线性）子流形上，该子流形可以通过后续的 ReduNet 层进行线性化并与其他类别分离，如稍后在 \Cref{fig:learn-to-classify-diagram} 中所示。

由这些多通道特征 $\bar\Z \doteq [\bar \z_1, \ldots, \bar \z_N] \in \Re^{C \times d \times N}$ 的循环版本（即 $\circm(\bar\Z) \doteq [ \circm(\bar\z_1), \dots, \circm(\bar\z_N)] \in \Re^{dC \times dN}$）构建的 ReduNet，保留了上述良好的不变性：线性算子（现在记为 $\bar{\bm E}$ 和 $\bar{\bm C}_k$）仍然是块循环的，并表示{\em 多通道一维循环卷积。} 具体而言，我们有以下结果。
\begin{proposition}[$\bar{\bm E}$ 和 $\bar{\bm C}_k$ 的多通道卷积结构]
    矩阵
    \begin{equation}
        \label{eq:def-E-bar}
        \bar{\bm E} \doteq \alpha\left(\bm I + \alpha\,
        \circm(\bar\Z) \circm(\bar\Z)^\top \right) ^{-1}
    \end{equation}
    是块循环的，即
    \begin{equation*}
        \bar{\bm E} =
        \left[
            \begin{matrix}
                \bar{\bm E}_{1, 1} & \cdots & \bar{\bm E}_{1, C}\\
                \vdots & \ddots & \vdots \\
                \bar{\bm E}_{C, 1} & \cdots & \bar{\bm E}_{C, C}\\
        \end{matrix}\right] \in \Re^{dC \times dC},
    \end{equation*}
    其中每个 $\bar{\bm E}_{c, c'}\in \Re^{d \times d}$ 都是一个循环矩阵。此外，$\bar{\bm E}$ 表示一个多通道循环卷积，即对于任何多通道信号 $\bar\z \in \Re^{C \times n}$，我们有
    $$\bar{\bm E} \cdot \textsf{vec}(\bar\z) = \textsf{vec}( \bar{\bm
    e} \circledast \bar\z).$$
    在上述公式中，$\bar{\bm e} \in \Re^{C \times C \times d}$ 是一个多通道卷积核，其中 $\bar{\bm e}[c, c'] \in \Re^{d}$ 是 $\bar{\bm E}_{c, c'}$ 的第一个列向量，而 $\bar{\bm e} \circledast \bar\z \in \Re^{C \times d}$ 是定义如下的多通道循环卷积：
    \begin{equation*}
        (\bar{\bm e} \circledast \bar\z)[c] \doteq \sum_{c'=1}^C
        \bar{\bm e}[c, c'] \circledast \bar{\z}[c'], \quad \forall c
        = 1, \ldots, C.
    \end{equation*}
    类似地，与 $\bar{\bm Z}$ 的任何子集相关联的矩阵 $\bar{\bm C}_k$ 也是多通道循环卷积。
    \label{prop:multichannel-circular-conv-1}
\end{proposition}
根据 \Cref{prop:multichannel-circular-conv-1}，平移不变的 ReduNet 在构造上是一个用于多通道一维信号的深度卷积网络。请注意，即使初始的提升核是分离的 \eqref{eqn:sparse-lifting}，在 \eqref{eq:def-E-bar} 中用于计算 $\bar{\bm E}$（对于 $\bar{\bm C_k}$ 也是如此）的矩阵求逆也会在所有 $C$ 个通道之间引入“串扰”。 % 因此，这些多通道卷积通常{\em 不是}深度可分离的，这与曾经被建议用于简化多通道卷积神经网络的 Xception 网络~\cite{Xception} 不同。\footnote{数据上的哪些额外结构会导致深度可分离卷积，这仍然是一个悬而未决的问题。}

\begin{remark}[在频域中降低计算复杂度]
    在 \eqref{eq:def-E-bar} 中计算 $\bar{\bm E}$ 需要对一个大小为 $dC \times dC$ 的矩阵求逆，这通常具有 $O(d^3C^3)$ 的复杂度。然而，利用循环矩阵可以通过离散傅里叶变换 (DFT) 矩阵对角化这一事实，可以显著降低复杂度。正如 \cite{chan2021redunet} 中所示，为了计算 $\bar{\bm E}$ 和 $\bar{\bm C}_k \in \Re^{dC \times dC}$，我们只需要在频域中计算 $d$ 次 $C\times C$ 块的逆，因此整体复杂度变为 $O(dC^3)$。

\end{remark}

\paragraph{整体网络架构与比较。}
根据上述推导，我们看到，为了找到多类信号/图像的平移不变的线性判别表示 (LDR)，稀疏编码、具有多通道卷积的多层架构、不同的非线性激活以及频谱计算，都成为有效且高效地实现该目标的{\em 必要}组件。\Cref{fig:learn-to-classify-diagram} 说明了通过对输入稀疏代码进行不变率衰减来学习这种表示的整体过程。

\begin{figure}[t]
    \centering
    \includegraphics[width=0.98\linewidth]{\toplevelprefix/chapters/chapter5/figs/learn_to_classify_diagram_updated.png}
    \caption{具有平移不变性的多类信号分类的整体过程：多通道提升、稀疏编码，随后是用于不变率衰减的多通道卷积 ReduNet。为了将平移不变的多类信号映射到不相干的（线性）子空间作为 LDR，这些组件是{\em 必要}的。请注意，大多数现代深度神经网络的架构都类似于这个过程。这样学习到的 LDR 有助于后续的分类等任务。}
    \label{fig:learn-to-classify-diagram}
\end{figure}

\begin{example}[数字的不变分类]
    接下来，我们提供 ReduNet 在真实的 10 类 MNIST 数据集上学习\textit{旋转}不变特征的经验性能。
    我们在图像 $\bm{x}\in\mathbb{R}^{H\times W}$ 上施加一个极坐标网格，其几何中心作为二维极坐标网格的中心（如 \Cref{fig:samples-invariant-1d-mnist-diagram} 所示）。
    对于每个半径 $r_i$，$i \in [C]$，我们可以针对每个角度 $\gamma_l =l\cdot({2\pi}/\Gamma)$ 采样 $\Gamma$ 个像素，其中 $l \in [\Gamma]$。
    然后，给定数据集中的样本图像 $\bm{x}$，我们将（采样后的）极坐标系中的图像表示为多通道信号 $\bm{x}_p \in\R^{\Gamma\times C}$。
    这里的目标是学习一个旋转不变的表示，即我们期望学习 $f(\cdot, \theta)$ 使得 $\{f(\bm{x}_p \circ \mathfrak{g}, \theta)\}_{\mathfrak{g} \in\mathbb{G}}$ 位于同一个子空间中，其中 $\mathfrak{g}$ 是极角的循环平移。
    我们使用 $N=100$ 个训练样本（每个类别 $10$ 个），并设置 $\Gamma=200$，$C=15$ 进行极坐标采样。
    通过在极坐标系中执行上述采样，我们可以获得数据矩阵 $\bm{X}_p \in \mathbb{R}^{(\Gamma\cdot C) \times N}$。
    对于 ReduNet，我们设置层数/迭代次数 $L=40$，精度 $\epsilon=0.1$，步长 $\eta=0.5$。在第一层之前，我们通过与 20 个大小为 5 的随机高斯核进行一维循环卷积来对输入进行提升。

    \begin{figure}[t]
        \begin{subfigure}[t]{0.3\textwidth}
            \centering
            \includegraphics[width=\textwidth]{\toplevelprefix/chapters/chapter5/figs/1d-rotation.pdf}
            \caption{$\bm{X}_{\text{rotation}}$}
        \end{subfigure}
        \hfill
        \begin{subfigure}[t]{0.3\textwidth}
            \centering
            \includegraphics[width=\textwidth]{\toplevelprefix/chapters/chapter5/figs/mnist1d_img0.pdf}
            \caption{$\bm{Z}_{\text{rotation}}$}
        \end{subfigure}
        \hfill
        \begin{subfigure}[t]{0.3\textwidth}
            \centering
            \includegraphics[width=\textwidth]{\toplevelprefix/chapters/chapter5/figs/mnist1d_img1.pdf}
            \caption{损失}
        \end{subfigure}
        \caption{\small MNIST 数字旋转图像的示例，每次旋转 18$^{\circ}$。(\textbf{左}) 极坐标表示示意图；(\textbf{右}) 数字 `0' 和 `1' 的旋转图像。}
        \label{fig:samples-invariant-1d-mnist-diagram}
    \end{figure}

    为了评估学习到的表示，每个训练样本都通过其 20 个旋转版本进行数据增强，每个版本的平移步长为 10。我们计算了 $m \times 20$ 个增强训练输入 $\bm{X}_{\text{rotation}}$ 之间的余弦相似度，结果如 \Cref{fig:redu-invariant-1d-mnist-diagram} (\textbf{a}) 所示。
    我们比较了所有增强版本学习到的特征（即 $\bar{\bm{Z}}_{\text{rotation}}$）之间的余弦相似度，并将结果总结在 \Cref{fig:redu-invariant-1d-mnist-diagram} (\textbf{b}) 中。
    正如我们所见，如此构建的旋转不变 ReduNet 能够将来自 10 个不同类别的训练数据（以及其所有旋转版本）映射到 10 个几乎正交的子空间中。也就是说，学习到的子空间对极角中的平移变换是真正不变的。接下来，我们随机抽取另外 $100$ 个测试样本，并遵循相同的增强过程。
    在 \Cref{fig:redu-invariant-1d-mnist-diagram} (\textbf{c}) 中，我们可视化了 ReduNet 在训练和测试数据集上第 $\ell$ 层表示的 MCR$^{2}$ 损失。
    从这些结果中，我们可以发现构建的 ReduNet 确实能够最大化 MCR$^{2}$ 损失，并能泛化到测试数据。

    \begin{figure}[t]
        \begin{subfigure}[t]{0.3\textwidth}
            \centering
            \includegraphics[width=\textwidth]{\toplevelprefix/chapters/chapter5/figs/mnist1d-heatmap-X_translate_train_all.pdf}
            \caption{$\bm{X}_{\text{rotation}}$}
        \end{subfigure}
        \hfill
        \begin{subfigure}[t]{0.3\textwidth}
            \centering
            \includegraphics[width=\textwidth]{\toplevelprefix/chapters/chapter5/figs/mnist1d-heatmap-Z_translate_train_all.pdf}
            \caption{$\bm{Z}_{\text{rotation}}$}
        \end{subfigure}
        \hfill
        \begin{subfigure}[t]{0.32\textwidth}
            \centering
            \includegraphics[width=\textwidth]{\toplevelprefix/chapters/chapter5/figs/mnist1d-loss-traintest.pdf}
            \caption{损失}
        \end{subfigure}
        \caption{\small (a)(b) 是旋转训练数据 $\bm{X}_{\text{rotation}}$ 和学习到的旋转不变特征 $\bar{\bm{Z}}_{\text{rotation}}$ 之间余弦相似度的热力图。(d) 可视化了跨层的训练/验证 MCR$^2$ 损失。}
        \label{fig:redu-invariant-1d-mnist-diagram}
    \end{figure}

\end{example}

\section{基于展开优化的白盒 Transformer}\label{sec:chap4-white-box-transformer}
正如我们在上一节中所看到的，我们利用分类问题为 ResNet 和 CNN 等流行深度网络的主要架构特征提供了严格的解释：此类网络的每一层都可以被视为模仿一个梯度步长，该步长增加了率衰减（或信息增益）目标。这种视角还引出了一个有些令人惊讶的事实：这种深度网络（即 ReduNet）各层的参数和算子可以以纯前向的方式计算。

尽管 ReduNet 在理论和概念上具有重要意义，但有几个因素限制了它的实用性。首先，如上所述，以前向方式计算每一层中矩阵算子的计算成本可能非常高。其次，这样计算出的算子在优化目标时可能不够有效，可能需要数千次迭代（因此需要数千层）。正如我们在 \Cref{sec:LISTA} 中关于 LISTA 所看到的，这两个问题可以通过允许优化这些算子并使它们通过反向传播变得可学习来解决。\footnote{或者，也许可以通过前向和反向优化的混合来实现。}

衍生出 ReduNet 的监督分类设置也有一定的局限性。在实践中，一张图像可能不属于单一类别，因为它可能包含多个对象。因此，假设图像的不同区域属于不同的低维模型（例如高斯分布或子空间）会更具一般性。正如我们将看到的，这种推广将导致一个既简单又通用的架构，它统一了我们在上一章中看到的率衰减和去噪操作。此外，由此获得的架构类似于流行的 Transformer 架构。

\subsection{稀疏率衰减的展开优化}

在过去的几年里，随着可用于训练深度网络的数据量和质量急剧增加，该领域已经从将输入视为“原子”转变为将输入视为“词元序列”，其中每个词元都是整体输入的一部分。这种视角的转变在哲学和实用性上都有所改进：在哲学上，它允许我们测量和建模输入的每个部分及其相互作用；在实用性上，它允许我们使用类似的深度网络（如 Transformer 等\textit{序列到序列模型}）来学习各种数据分布。现在，我们需要理解在将数据视为词元序列的“以词元为中心”的模型中的表征学习，并应对上述挑战（例如，在率衰减框架的背景下，没有每个词元的标签）。

形式上，令 \(\X = \mat{\x_{1}, \dots, \x_{N}} \in \bR^{D \times N}\) 表示代表我们数据源的随机变量。在视觉任务中，每个 \(\x_{i} \in \bR^{D}\) 被解释为一个\textit{词元}，通常对应于一个图像块。在语言任务中，每个 \(\x_{i} \in \bR^{D}\) 被解释为一个\textit{词元嵌入}，即离散词元（如单词或子词）的连续向量表示。\footnote{为了方便起见，我们在本章中稍微滥用一下术语，将离散词元及其相关的嵌入统称为词元。}  %
\(\x_{i}\) 可能具有任意的相关结构。我们使用 \(\Z = \mat{\z_{1}, \dots, \z_{N}} \in \bR^{d \times N}\) 来表示定义我们表示的随机变量，其中 \(\z_{i} \in \bR^{d}\) 是对应词元 \(\x_i \in \bR^{D}\) 的表示。%

\begin{remark}
    在 Transformer 中，每个输入样本通常被转换为一个{\em 词元}序列。词元是从原始输入中提取的基本信息单元：在自然语言处理中，词元通常是单词或子词；在计算机视觉中，它们对应于图像块；在其他模态中，它们可能表示时间步长、空间位置或其他特定领域的单元。{\em 词元嵌入}是词元的连续向量表示，作为 Transformer 的输入。它将每个词元映射到高维空间中的一个点，使模型能够使用数值计算来处理符号输入。
    {\em 词元表示}是编码词元语义或结构信息的向量，通常由 Transformer 的中间层或最终层产生。这些表示旨在捕获对下游任务（如分类、生成或回归）有用的输入的有意义特征。有关此词汇表及其与实现的关系的更多详细信息，请参阅 \Cref{sec:contrastive_learning}。
\end{remark}

\paragraph{学习结构化和紧凑表示的目标。}
受前面关于率衰减讨论的启发（\Cref{sec:chap4-white-box-model-via-unrolling}），我们提出，在这种以词元为中心的数据模型中，表征学习的目标是找到一个特征映射 \(f \colon \X \in \bR^{D \times N} \to \Z\in \bR^{d \times N}\)，它将具有潜在非线性和多模态分布的输入词元 \(\{\bm x_i\}_{i=1}^N \subset \R^D\) 转换为（分段）\textit{线性化且紧凑}的词元表示 \(\{\bm z_i\}_{i=1}^N \subset \bR^{d}\)。虽然词元表示 \(\{\z_{i}\}_{i = 1}^{N}\) 的{\em 联合}分布可能很复杂（且特定于任务），但我们进一步认为，要求单个词元表示的目标{\em 边缘}分布高度压缩和结构化、易于紧凑编码，是合理且实用的。特别是，我们要求该分布是\textit{低维（例如 \(K\) 个）高斯分布的混合}，使得第 \(k\) 个高斯分布的均值为 \(\Zero \in \bR^{d}\)，协方差为 \(\vSigma_{k} \succeq \Zero \in \bR^{d \times d}\)，并且其支撑集由正交基 \(\vU_{k} \in \bR^{d \times p}\) 张成。
我们将所有高斯分布的基的集合记为 $\vU_{[K]} = \{\vU_k\}_{k=1}^K$。因此，为了最大化最终词元表示的\textit{信息增益} \cite{ma2022principles}，我们希望最大化它们的率衰减（见 \Cref{subsec:MCR2}），即：
\begin{align}\label{eq:rate reduction}
    \mathrm{max}_{\bm Z \in \R^{d\times N}}\ \Delta R_{\epsilon}(\Z
    \mid \vU_{[K]}) \doteq R_{\epsilon}(\Z) - R^c_{\epsilon}(\Z \mid \vU_{[K]}).
\end{align}
这里，第一项 $R_{\epsilon}$ 是对整个词元表示集合的有损编码率的估计。更具体地说，如果我们将词元表示 $\{\bm z_i\}_{i=1}^N$ 视为来自单个零均值高斯分布的独立同分布 (i.i.d.) 样本，则在量化精度 $\epsilon > 0$ 下，它们的有损编码率由下式给出：
\begin{equation}\label{eq:coding_rate}
    R_{\epsilon}(\Z) \doteq \frac{1}{2}\textrm{logdet}\left(\I +
    \frac{d}{N\epsilon^{2}}\Z^{\top}\Z\right) =
    \frac{1}{2}\textrm{logdet}\left(\I +
    \frac{d}{N\epsilon^{2}}\Z\Z^{\top}\right).
\end{equation}
第二项 $R_{\epsilon}^c$ 是在码本 $\bm U_{[K]}$ 下的有损编码率的估计，由下式给出：
\begin{equation}
    \begin{aligned}
        R_{\epsilon}^c(\Z \mid \vU_{[K]}) &\doteq
        \sum_{k=1}^{K}R_{\epsilon}(\vU_k^{\top} \Z) =
        \frac{1}{2}\sum_{k=1}^{K}\log\det\left(\I +
        \frac{p}{N\epsilon^{2}}(\vU_k^{\top}\Z)^{\top}(\vU_k^{\top}\Z)\right).
    \end{aligned}
    \label{eq:def-mcr-Rc}
\end{equation}

\begin{remark}
    编码率的表达式 \eqref{eq:def-mcr-Rc} 可以看作是原始率衰减目标 \eqref{eq:MCRc} 中使用的编码率 \(R_\epsilon^{c}\) 的推广。特别是，原始目标是针对特定数据实现 \(\vX\) 的一组已知成员标签 \(\{\vPi_{k}\}\) 定义的。相比之下，当前目标是针对子空间 \(\vU_{[K]}\) 定义的，这些子空间独立于任何特定的实现，但被假设为支撑词元表示的分布。
    假设一个词元表示 $\bm z_i$ 属于子空间 $\bm U_k$，并且这些子空间彼此近似正交，即对于所有 $k \neq l$，$\bm U_k^\top \bm U_l \approx \bm 0$。那么，可以验证投影 $\bm U_k\bm U_k^\top \bm z_i  = \bm z_i$ 且对于所有 $l \neq k$，$\bm U_l\bm U_l^\top \bm z_i \approx \bm 0$。这些正交投影有效地充当了隐式成员标签，标识了每个词元表示所属的子空间。
\end{remark}

\begin{figure}[t!]
    \centering
    \includegraphics[width=0.95\textwidth]{\toplevelprefix/chapters/chapter5/figs/coding-transform.png}
    \caption{ \small\textbf{通过率衰减和稀疏性对三组表示的比较。} 每个 $S_i$ 代表一个线性子空间，蓝球的数量代表编码率之差 $\Delta R_{\epsilon}(\vZ \mid \vU_{[K]}) = R_\epsilon(\vZ) - R^c_\epsilon(\vZ \mid \vU_{[K]})$。
    }
    \label{fig:sparse-rate-reduction-diagram}
\end{figure}

\paragraph{稀疏率衰减。} 请注意，率衰减目标 \eqref{eq:rate reduction} 对表示和子空间的任意联合旋转是不变的。特别是，优化率衰减目标可能不会自然地导致轴对齐（即\textit{稀疏}）的表示。{例如，考虑 \Cref{fig:sparse-rate-reduction-diagram} 中学习到的三组表示。编码率衰减从 (a) 到 (b) 增加，但由于它在旋转下是不变的，因此从 (b) 到 (c) 保持不变。} 因此，我们希望对表示（及其支撑子空间）进行变换，使得表示 $\vZ$ 最终相对于所得表示空间的标准坐标变得稀疏\footnote{具体而言，具有很少的非零项。}{如 \Cref{fig:sparse-rate-reduction-diagram}(c) 所示}。因此，为了确保最终表示易于进行更紧凑的编码，我们希望对表示（及其支撑子空间）进行变换，使它们相对于所得表示空间的标准坐标变得\textit{稀疏}。\footnote{也就是说，具有最少的非零项。}
在计算上，我们可以将上述两个目标结合成一个统一的优化目标：
\begin{equation}
    \max_{f \in \mathcal{F}}\ [\Delta R_{\epsilon}(\Z \mid \vU_{[K]})
    - \lambda \|\bm{Z}\|_0] \qquad \text{s.t.}\ \Z = f(\X),
    \label{eq:sparse-rr}
\end{equation}
其中 $\mathcal{F}$ 表示一般的函数类，$\ell_0$ 范数 $\|\Z\|_0$ 促进了最终词元表示 \(\Z = f(\X)\) 的稀疏性。%

在实践中，$\ell_0$ 范数通常被松弛为 $\ell_1$ 范数，以提高计算的可追踪性并启用凸优化技术 \cite{Wright-Ma-2022}。受此启发，我们相应地松弛了问题 \eqref{eq:sparse-rr}，从而得到一个既忠实于原始稀疏性目标，又更易于使用高效算法的公式，如下所示：
\begin{equation}
    \begin{aligned}
        \max_{f \in \mathcal{F}}\ [\Delta R_{\epsilon}(\Z \mid
        \vU_{[K]}) - \lambda \|\bm{Z}\|_1]  \qquad \text{s.t.}\ \Z = f(\X),
        \label{eq:sparse-rr-1}
    \end{aligned}
\end{equation}
稍微滥用一下术语，我们通常也将这个目标函数称为\textit{稀疏率衰减}。

\paragraph{通过展开优化构建白盒网络架构。}

尽管陈述起来很简单，但上述目标中的每一项在计算上都具有优化的挑战性 \cite{Wright-Ma-2022}。因此，很自然地采用一种近似方法，通过级联多个（例如 $L$ 个）简单的\textit{增量和局部}操作 $f^\ell$ 来实现全局变换 $f$ 以优化 \eqref{eq:sparse-rr}，这些操作将表示分布推向所需的简约模型分布：
\begin{equation}
    f\colon \X \xrightarrow{\hspace{1mm} f^{\mathrm{pre}} \hspace{1mm}}
    \Z^1 \rightarrow \cdots \rightarrow \Z^\ell
    \xrightarrow{\hspace{1mm} f^{\ell} \hspace{1mm}} \Z^{\ell+1}
    \rightarrow  \cdots \xrightarrow{\hspace{1mm} f^{L}} \Z^{L+1} = \Z,
    \label{eq:incremental}
\end{equation}
其中 $f^{\mathrm{pre}}: \bR^{D} \rightarrow \bR^{d}$ 是预处理映射，它将每个输入词元 $\x_{i} \in \bR^{D}$ 转换为初始词元表示 $\z_{i}^{1} \in \bR^{d}$。
每个增量\textit{前向映射} $\Z^{\ell + 1} = f^\ell(\Z^\ell)$（或称为“层”），在以其输入 $\Z^\ell$ 的分布为条件的情况下，转换词元分布以\textit{优化}上述稀疏率衰减目标 \eqref{eq:sparse-rr}。

\begin{remark}
    与 ReduNet 等其他展开优化方法（见 \Cref{sec:chap4-white-box-model-via-unrolling}）相比，我们在每一层\textit{显式地建模} $\Z^\ell$ 的分布，例如将其建模为线性子空间的混合或由字典稀疏生成。模型参数是从数据中学习的（例如通过端到端训练的\textit{反向传播}）。这种前向“优化”和反向“学习”之间的分离，阐明了每一层作为转换其输入分布的算子的数学作用，而输入分布反过来又由该层的参数建模（并随后被学习）。
\end{remark}

现在，我们展示如何通过展开优化的视角推导这些增量和局部操作，以解决问题 \eqref{eq:sparse-rr-1}。一旦我们决定使用增量方法来优化问题 \eqref{eq:sparse-rr-1}，就有多种可能的选择来实现优化。给定 $\Z^\ell$ 的模型（例如子空间的混合 $\vU_{[K]}$），我们选择一种具有坚实概念基础的两步\textit{交替最小化}方法。首先，我们通过梯度下降\textit{压缩}词元 $\vZ^{\ell}$，以最小化编码率项 $R^c_\epsilon (\vZ \mid \vU_{[K]}^\ell)$。具体来说，我们以学习率 $\kappa$ 对 $R^c_\epsilon$ 采取梯度步长，如下所示：
\begin{align}\label{eq:Z l+1/2}
    \bm Z^{\ell+1/2} = \bm Z^\ell - \kappa \nabla_{\bm Z}
    R^c_\epsilon (\vZ \mid \vU_{[K]}^\ell).
\end{align}
接下来，我们对压缩后的词元进行\textit{稀疏化}，通过适当松弛的近端梯度步长生成 \(\vZ^{\ell + 1}\)，以最小化剩余项 $\lambda \norm{\vZ}_{1} - R_{\epsilon}(\vZ)$。正如我们稍后将详细论证的那样，我们可以通过求解关于字典 $\bm D^\ell$ 的稀疏表示问题来找到这样的 $\bm Z^{\ell+1}$：
\begin{equation}\label{eq:Z l+1}
    \vZ^{\ell+1} = \argmin_{{\vZ}}  \bigg\{\lambda \norm{\vZ}_1 +
    \frac{1}{2}\norm{\vZ^{\ell + 1/2} - \vD^{\ell} {\vZ}}_F^2\bigg\}.
\end{equation}
在下文中，我们提供了上述两个步骤中每一步的技术细节，并推导了用于其实现的高效更新。

\paragraph{自注意力作为词元表示编码率的梯度下降。} 对于第一步 \eqref{eq:Z l+1/2}，编码率的梯度 \(\nabla_{\bm Z} R^c_\epsilon\) 的计算成本很高，因为它涉及 \(K\) 个独立的矩阵求逆，每个求逆对应于具有基 \(\vU_{k}^{\ell}\) 的 \(K\) 个子空间之一：
\begin{equation}
    \nabla_{\bm Z} R_{\epsilon}^c(\vZ \mid \vU_{[K]})
    = \frac{p}{N\epsilon^2}\sum_{k=1}^K \vU_k\vU_k^\top\vZ\Big(\I +
    \frac{p}{N\epsilon^2}(\vU_k^\top\vZ)^\top(\vU_k^\top\vZ)\Big)^{-1}.
    \label{eq:rate-gradient}
\end{equation}
现在，我们证明这个梯度可以自然地使用所谓的“多头子空间自注意力” (MSSA) 算子来近似，该算子具有与具有 \(K\) 个头（即每个子空间一个头，来自每个矩阵求逆）的多头自注意力算子 \citep{vaswani2017attention} 相似的函数形式。在这里，我们使用一阶诺伊曼级数（见 \Cref{ex:neumannn}）来近似梯度 \eqref{eq:rate-gradient}：
\begin{align}\label{eq:gd_rc_neumann}
    \nabla_{\bm Z} R_{\epsilon}^{c}(\vZ \mid \vU_{[K]})
    &\approx \frac{p}{N\epsilon^2} \sum_{k =
    1}^{K}\vU_{k}\vU_{k}^\top\vZ\left(\vI - \frac{p}{N\epsilon^2}
    (\vU_{k}^\top\vZ)^\top(\vU_{k}^\top\vZ)\right) \notag \\
    &{= \frac{p}{N\epsilon^2} \left(\sum_{k = 1}^{K}
        \vU_{k}\vU_{k}^\top\right)\vZ -  \left(
        \frac{p}{N\epsilon^2}\right)^2\sum_{k = 1}^{K}
    \vU_{k}(\vU_{k}^\top\vZ)(\vU_{k}^\top\vZ)^\top(\vU_{k}^\top\vZ)}.
\end{align}
在这种近似中，我们通过投影特征之间的自相关 $(\vU_{k}^\top\vZ)^\top(\vU_{k}^\top\vZ)$ 来计算投影词元表示 $\{\bm U_k^\top\bm z_i\}$ 之间的相似度，并使用 softmax 将其转换为成员分布，即 $\softmax{(\vU_{k}^\top\vZ)^\top(\vU_{k}^\top\vZ)}$。
假设子空间的并集 $\bm U_{[K]}$ 张成了整个空间。那么，我们有 $\sum_{k = 1}^{K} \vU_{k}\vU_{k}^\top = \bm I$。因此，\eqref{eq:gd_rc_neumann} 变为：
\begin{align}\label{eq:grad Rc}
    \nabla_{\bm Z} R_{\epsilon}^{c}(\vZ \mid \vU_{[K]})
    \approx  \frac{p}{N\epsilon^2} \vZ -  \left(
    \frac{p}{N\epsilon^2}\right)^2 \MSSA\left(\vZ^{\ell} \mid
    \vU_{[K]}^{\ell}\right),
\end{align}
其中 MSSA 通过 SSA 算子定义如下：
\begin{align}
    & \mathrm{SSA}\left(\vZ \mid \vU_{k}\right)
    \doteq (\vU_{k}^\top
    \vZ)\mathrm{softmax}\left((\vU_{k}^\top\vZ)^\top(\vU_{k}^\top\vZ)\right),
    \ \forall k \in [K], \label{eq:SSA} \\
    & \mathrm{MSSA}\left(\vZ \mid \vU_{[K]}\right)
    \doteq \frac{p}{N\epsilon^2} \mat{\vU_{1}, \dots,
    \vU_{K}}\mat{\mathrm{SSA}({\vZ \mid \vU_{1}}) \\ \vdots
    \\ \mathrm{SSA}({\vZ \mid \vU_{K}})}.\label{eq:Multi-Head-SSA}
\end{align}
将 \eqref{eq:grad Rc} 代入 \eqref{eq:Z l+1/2} 可得，它可以自然地近似为：
\begin{equation}
    \vZ^{\ell + 1/2} = \left(1 - \frac{\kappa p}{N\epsilon^2}\right)
    \vZ^{\ell} + \frac{\kappa p}{N\epsilon^2}
    \mathrm{MSSA}\left(\vZ^{\ell}\ \middle|\ \vU_{[K]}^{\ell}\right).
    \label{eq:gd-mcr-parts}
\end{equation}

\begin{remark}
    \eqref{eq:SSA} 中的 SSA 算子类似于典型 Transformer \citep{vaswani2017attention} 中的\textit{注意力算子}，不同之处在于，这里的值、键和查询的线性算子都被设置为与子空间基\textit{相同}，即 $\vV_{k} = \vK_{k} = \bm{Q}_{k} = \vU_k^*$。因此，我们将 $\mathrm{SSA}({\spcdot\mid\vU_k}): \bR^{d\times n} \rightarrow \bR^{p\times n}$ 命名为\textbf{子空间自注意力} (SSA) 算子。然后，\eqref{eq:Multi-Head-SSA} 中的整个 MSSA 算子，形式上定义为 \(\mathrm{MSSA}({\spcdot \mid \vU_{[K]}}) \colon \bR^{d \times n} \to \bR^{d \times n}\)，被称为\textbf{多头子空间自注意力} (MSSA) 算子，它通过使用依赖于模型的权重进行平均来聚合注意力头的输出，这在概念上类似于现有 Transformer 网络中流行的多头自注意力算子。整体梯度步长 \eqref{eq:gd-mcr-parts} 类似于 Transformer 中带有跳跃连接的多头自注意力。
\end{remark}

\paragraph{MLP 作为词元表示稀疏编码的近端梯度下降。} 对于交替最小化的第二步，我们需要最小化 $\lambda \norm{\vZ}_{1} - R_{\epsilon}(\vZ)$。请注意，梯度 \(\nabla R_{\epsilon}(\vZ)\) 涉及矩阵求逆，因此使用朴素的近端梯度（见 \Cref{subsec:pgd}）来优化此问题在大规模问题上变得难以处理。 %
因此，我们采用一种不同的、简化的方法来在表示多样性和稀疏化之间进行权衡：我们假设一个（完备的）不相干或正交字典 $\vD^{\ell} \in \bR^{d \times d}$，并要求相对于 \(\vD^{\ell}\) 对中间迭代 $\vZ^{\ell + 1/2}$ 进行稀疏化。也就是说，$\vZ^{\ell + 1/2} \approx \vD^{\ell} \vZ^{\ell + 1}$，其中 $\vZ^{\ell + 1}$ 更稀疏；也就是说，它是 \(\vZ^{\ell + 1/2}\) 的\textit{稀疏编码}。字典 \(\vD^{\ell}\) 用于同时稀疏化所有词元。
根据不相干性假设，我们有 $(\vD^{\ell})^\top(\vD^{\ell}) \approx \vI$。因此，由 \eqref{eq:coding_rate} 我们有：
\begin{equation}
    R_{\epsilon}(\vZ^{\ell + 1/2}) \approx
    R_{\epsilon}(\vD^{\ell}\vZ^{\ell + 1}) \approx R_{\epsilon}(\vZ^{\ell + 1}).
\end{equation}
为了求解 $\lambda \norm{\vZ}_{1} - R_{\epsilon}(\vZ)$，我们优化以下问题：
\begin{align*}
    \vZ^{\ell + 1} \approx \argmin_{\vZ}  \norm{\vZ}_1 \quad
    \mbox{subject to} \quad \vZ^{\ell + 1/2} = \vD^{\ell}\vZ.
\end{align*}
{上述稀疏表示规划通常通过将其松弛为无约束凸规划来求解，即 LASSO \citep{Wright-Ma-2022}：}
\begin{equation}
    \vZ^{\ell + 1} \approx \argmin_{\vZ} \left[\lambda \norm{\vZ}_1 +
    \frac{1}{2}\norm{\vZ^{\ell + 1/2} - \vD^{\ell} \vZ}_F^2 \right].
\end{equation}
在我们的实现中，我们还对 $\vZ^{\ell + 1}$ 添加了非负约束，并求解相应的非负 LASSO：
\begin{equation}
    \vZ^{\ell + 1} \approx \argmin_{\vZ \geq \bm 0}
    \left[\lambda\norm{\vZ}_1 + \frac{1}{2}\norm{\vZ^{\ell + 1/2} -
    \vD^{\ell} \vZ}_{F}^{2}\right].
    \label{eq:sparse-nonnegative}
\end{equation}
然后，我们通过执行展开的{\em 近端梯度下降}步长（称为 ISTA 步长 \citep{beck2009fast}）来增量优化 \Cref{eq:sparse-nonnegative}，从而给出更新：
\begin{align}\label{eq:ista-block}
    \vZ^{\ell + 1}
    &= \mathrm{ISTA}({\vZ^{\ell + 1/2} \mid \vD^{\ell}}), \\
    \text{where} \quad \mathrm{ISTA}({\vZ \mid \vD})
    &\doteq \operatorname{ReLU}(\vZ - \eta \vD^\top(\vD\vZ - \vZ) -
    \eta \lambda \bm{1}).
\end{align}

\begin{figure}
    \centering
    \includegraphics[width=0.75\textwidth]{\toplevelprefix/chapters/chapter5/figs/crate_encoder_architecture.pdf}
    \caption{\small \textbf{CRATE 编码器架构的一层。} 完整的架构只是这些层的级联，并带有一些初始的分词器、预处理头以及最终的特定任务头（即分类头）。}
    \label{fig:crate_backbone}
\end{figure}

\subsection{整体白盒 Transformer 架构：CRATE}

我们现在通过展开上述更新来设计一个白盒 Transformer 架构，命名为编码率 Transformer (\textsc{crate})。通过结合上述两个步骤 \eqref{eq:gd-mcr-parts} 和 \eqref{eq:ista-block}：
\begin{enumerate}
    \item 将样本内的词元局部压缩为子空间混合结构，从而产生多头子空间自注意力块 —— \texttt{MSSA}；
    \item 通过稀疏编码对所有样本的词元集进行全局稀疏化，从而产生稀疏化块 —— \texttt{ISTA}；
\end{enumerate}
我们可以得到以下基于率衰减的 Transformer 层，如 \Cref{fig:crate_backbone} 所示：
\begin{equation}
    \vZ^{\ell+1/2} \doteq \vZ^{\ell} + \texttt{MSSA}(\vZ^{\ell} \mid
    \vU_{[K]}^{\ell}),
    \qquad
    \vZ^{\ell+1}\doteq \texttt{ISTA}(\vZ^{\ell+1/2} \mid \bm D^\ell).
\end{equation}
按照 \eqref{eq:incremental} 中表示的增量构造组合多个这样的层，我们获得了一个白盒 Transformer 架构，该架构将数据词元转换为紧凑且稀疏的不相干子空间并集，其中 $f^{\pre}: \bR^{D \times N} \rightarrow \bR^{d \times N}$ 是预处理映射，它将输入词元 $\vX \in \bR^{D \times N}$ 转换为第一层表示 $\vZ^{1} \in \bR^{d \times N}$。该架构的整体流程如 \Cref{fig:crate-diagram} 所示。

\begin{figure}[t!]
    \centering
    \includegraphics[width=\textwidth]{\toplevelprefix/chapters/chapter5/figs/CRATE_fig1_patches.pdf}
    \vspace{-0.1in}
    \caption{
        \textbf{\textsc{crate} 白盒深度网络设计的“主循环”。}
        在将输入数据编码为词元序列 $\vZ^0$ 之后，\textsc{crate} 构建了一个深度网络，该网络通过针对分布的局部模型进行连续的{\textit{\textbf{压缩}}}（生成 $\vZ^{\ell+1/2}$），以及针对全局字典进行{\textit{\textbf{稀疏化}}}（生成 $\vZ^{\ell+1}$），将数据转换为低维子空间的规范配置。重复堆叠这些块并通过反向传播训练模型参数，可以产生强大且可解释的数据表示。
    }
    \label{fig:crate-diagram}
\end{figure}

\begin{remark}[\textbf{前向传递和反向传播的作用}]\label{sub:forward_backward}
    与 ReduNet \cite{chan2021redunet} 等其他展开优化方法相比，我们在每一层\textit{显式地建模}每个 $\vZ^\ell$ 和 $\vZ^{\ell + 1/2}$ 的分布，要么通过线性子空间的混合，要么通过字典稀疏生成。我们引入了这样一种解释：在每一层 \(\ell\)，学习到的子空间基 \(\vU_{[K]}^{\ell}\) 和学习到的字典 \(\vD^{\ell}\) 共同作为一个\textit{码本}或\textit{分析滤波器}，对每一层 \(\ell\) 的中间表示进行编码和转换。由于第 \(\ell\) 层的输入分布首先由 \(\vU_{[K]}^{\ell}\) 建模，然后由 \(\vD^{\ell}\) 转换，因此每一层的输入分布是不同的，所以我们需要在每一层有一个单独的码本以获得最简约的编码。这些码本的参数（即子空间基和字典）迄今为止被假设为完全已知，实际上是从数据中学习的（例如通过端到端训练中的\textit{反向传播}）。

    因此，我们的方法在由此推导出的白盒深度神经网络的前向“优化”和反向“学习”之间具有清晰的概念分离。也就是说，在其前向传递中，我们将每一层解释为一个算子，该算子以其输入分布的学习模型（即码本）为条件，将该分布转换为更简约的表示。在其反向传播中，该模型针对每一层输入分布的码本被更新，以更好地拟合某种（监督的）输入输出关系，如 \Cref{fig:forward-backward} 所示。这种概念上的解释意味着模型表示对特定任务和损失具有一定程度的不可知性；特别是，许多类型的任务和损失将确保每一层的模型得到拟合，从而确保模型产生简约的表示。
\end{remark}

我们现在通过测量所提出的网络 \textsc{crate} 在 ImageNet-1K 上的 top-1 分类准确率以及在几个广泛使用的下游数据集上的迁移学习性能，来展示其经验性能。
我们将结果总结在 \Cref{tab:crate_comparison_with_sota} 中。
迁移学习的方法是从预训练网络初始化，使用交叉熵损失进行微调。
由于设计的白盒 Transformer 架构在注意力块 (\texttt{MSSA}) 和非线性块 (\texttt{ISTA}) 中都利用了参数共享，因此 \textsc{crate}{-Base} 模型（2280 万参数）的参数数量与 ViT-Small（2205 万参数）~\cite{dosovitskiy2020image} 相似，并且不到相同配置的 ViT-Base（8654 万参数）参数数量的 30\%。
从 \Cref{tab:crate_comparison_with_sota} 中，我们发现，在模型参数数量相似的情况下，我们提出的网络实现了与 ViT 相似的 ImageNet-1K 和迁移学习性能，同时具有简单且有原则的设计。此外，在相同的训练超参数集下，我们在 \textsc{crate} 中观察到了有希望的缩放行为——我们通过扩大模型规模持续提高了性能。总而言之，\textsc{crate} 通过直接实现我们有原则的架构，在现实世界的大规模数据集上取得了有希望的性能。我们将在最终的应用 \Cref{ch:applications} 中提供有关图像分类实验结果的实现和分析的更多细节。

\begin{figure}[t]
    \begin{subfigure}[t]{0.48\textwidth}
        \centering
        \includegraphics[width=\textwidth]{\toplevelprefix/chapters/chapter5/figs/forward.png}
        \caption{\bf 前向传递}
    \end{subfigure}
    \hfill
    \begin{subfigure}[t]{0.48\textwidth}
        \centering
        \includegraphics[width=\textwidth]{\toplevelprefix/chapters/chapter5/figs/backward.png}
        \caption{\bf 反向传播}
    \end{subfigure}
    \caption{\small {\bf 深度网络中前向传递和反向传播的作用}。(a) 给定固定的子空间和字典 $\{(\bm U_{[K]}^{\ell}, \bm D^{\ell})\}_{\ell=1}^L$，每一层在前向传递中对表示执行压缩和稀疏化；(b) 反向传播从训练数据中学习子空间和字典 $\{(\bm U_{[K]}^{\ell}, \bm D^{\ell})\}_{\ell=1}^L$。 }
    \label{fig:forward-backward}
\end{figure}

\begin{table*}[t!]
    \centering
    \caption{\small 在 ImageNet-1K 上预训练时，\textsc{crate} 在具有不同模型规模的各种数据集上的 Top-1 分类准确率。对于 ImageNet-1K/ImageNet-1K ReaL，我们直接评估 top-1 准确率。对于其他数据集，我们使用在 ImageNet 上预训练的模型作为初始化，并通过微调评估迁移学习性能。}
    \label{tab:crate_comparison_with_sota}
    \small
    \setlength{\tabcolsep}{13.6pt}
    \resizebox{0.98\textwidth}{!}{
        \begin{tabular}{@{}lcccc|cc@{}}
            \toprule
            \textbf{模型} & \textsc{crate}{-T}  &
            \textsc{crate}{-S} & \textsc{crate}{-B} &
            \textsc{crate}{-L} & { \color{gray} ViT-T} &  {
            \color{gray}ViT-S } \\
            \midrule
            \midrule
            参数量 & 6.09M & 13.12M & 22.80M & 77.64M & {
            \color{gray} 5.72M} & { \color{gray} 22.05M} \\
            \midrule
            ImageNet-1K & 66.7 & 69.2 & 70.8 & 71.3 & { \color{gray}
            71.5} & { \color{gray} 72.4} \\
            ImageNet-1K ReaL & 74.0 & 76.0 & 76.5 & 77.4 & {
            \color{gray} 78.3 } & { \color{gray} 78.4} \\
            \midrule
            CIFAR10 & 95.5 & 96.0 & 96.8 & 97.2 & { \color{gray}
            96.6} & { \color{gray} 97.2} \\
            CIFAR100 & 78.9 & 81.0 & 82.7 & 83.6 & { \color{gray}
            81.8} & { \color{gray} 83.2}\\
            Oxford Flowers-102 & 84.6 & 87.1 & 88.7 & 88.3 & {
            \color{gray} 85.1} & { \color{gray} 88.5}\\
            Oxford-IIIT-Pets & 81.4 & 84.9 & 85.3 & 87.4 & {
            \color{gray} 88.5} & { \color{gray} 88.6} \\
            \bottomrule
    \end{tabular}}
\end{table*}

\section{通过设计得到的深度架构变体}
\label{sec:chap4-derive-white-box-transformer-variants}

到目前为止，我们希望我们已经提供了令人信服的证据，证明（流行的）深度网络的作用是实现某些优化算法，以最小化学习表示的编码率（或最大化信息增益）。然而，熟悉优化方法的读者可能已经注意到，上述架构（ReduNet 或 CRATE）对应于相当基础的优化技术。它们在效率或有效性方面可能有很大的改进空间。此外，如果我们相信所提出的用于解释深度网络的理论框架是正确的，那么它不仅应该有助于解释现有的架构，还应该指导我们开发更高效、更有效的架构。在本节中，我们展示了情况可能确实如此：由此产生的新架构不仅完全可解释，而且具有保证的正确性和提高的效率。

\subsection{纯注意力 Transformer 架构} \label{sub:aot}

在本小节中，我们提出了一种极简的 Transformer 架构，该架构由基于 MSSA 算子的可解释层组成。为了推导出一个仅包含必要组件的完全可解释的 Transformer 架构，我们认为表征学习的目标是将一组嘈杂的初始词元表示压缩为低维子空间的混合。 %
在这里，我们假设初始词元表示 $\bm Z^{(1)}$ 是从受噪声扰动的低秩高斯混合中采样的，如下所示：
\begin{definition}\label{def:MoG}
    令 $C_1,\dots,C_K$ 为索引集 $[N]$ 的一个划分，$\bm U_k \in \mathcal{O}^{d \times p_k}$ 表示每个 $k \in [K]$ 的第 $k$ 个子空间的正交基。如果对于每个 $k \in [K]$，
    \begin{align}\label{eq:MoG}
        \bm z_i = \underbrace{\bm U_{k} \bm a_i}_{\bf 信号} +
        \underbrace{\sum_{j \neq k}^K \bm U_j \bm e_{i,j}}_{\bf
        噪声},\ \forall i \in C_k,
    \end{align}
    其中对于所有 $i \in C_k$ 和 $k \in [K]$，$\bm{a}_i \overset{i.i.d.}{\sim} \mathcal{N}(\bm{0},\bm{I}_{p_k})$ 且 $\bm{e}_{i,j} \overset{i.i.d.}{\sim} \mathcal{N}(\bm{0},\delta^2\bm{I}_{p_j})$，$\{\bm{a}_i\}$ 和 $\{\bm{e}_{i,j}\}$ 分别相互独立，且 $\{\bm{a}_i\}$ 独立于 $\{\bm{e}_{i,j}\}$，则我们称词元表示 $\{\bm z_i \}_{i=1}^N \subseteq \R^d$ 是从嘈杂的低秩高斯混合分布中采样的。
\end{definition}
该模型作为近似现实世界预训练 LLM 中词元表示的理想化框架。它假设词元表示是从带有噪声的多个低秩高斯分布的混合中采样的。在该模型下，表征学习的目标是将一组嘈杂的初始词元表示压缩到相应的子空间中。此外，该模型与关于预训练大型语言模型中词元表示结构的两个公认假设非常吻合：“线性表示假设” \citep{jiang2024origins,park2023linear} 和“叠加假设” \citep{elhage2022toy,yun2021transformer}。

\begin{remark}
    线性表示假设认为，LLM 中的词元表示位于编码语义特征的低维线性子空间中。类似地，叠加假设表明，这些表示可以近似地表达为这些特征向量的稀疏线性组合。在 \Cref{def:MoG} 中，子空间的每个基 $\bm U_k$ 都可以解释为一组语义特征，其中每个特征对应于词元含义的特定方面。然后，词元表示被近似地表达为这些子空间基的稀疏线性组合，捕获词元的基本语义成分，同时忽略不相关的维度。
\end{remark}

\paragraph{词元表示的去噪算子。} 现在，我们展示 MSSA 算子（见 \eqref{eq:Multi-Head-SSA}）可以对由上述模型生成的词元表示进行增量去噪。具体来说，我们考虑对于每个 $\ell =1 ,\dots,L$，
\begin{align}\label{eq:MSSA}
    \bm Z^{\ell+1} =  \bm Z^{\ell} + \eta \sum_{k=1}^K \bm U_k\bm
    U_k^T \bm Z^{\ell} \varphi \left(\bm Z^{\ell^T}\bm U_k\bm
    U_k^T\bm Z^{\ell} \right),
\end{align}
其中 $\{\bm U_k\}_{k=1}^K$ 在 \Cref{def:MoG} 中定义，$\eta > 0$ 是步长，$\varphi$ 是逐元素算子，例如 softmax、ReLU 或其他函数。为了简化我们的推导，我们假设 \Cref{def:MoG} 中的子空间彼此正交，即对于所有 $k \neq j$，$\bm U_k^T\bm U_j = \bm 0$。请注意，这个假设并不具有限制性，因为在高维空间中，随机的低维子空间以高概率彼此不相干，即 $\bm U_k^T\bm U_j \approx \bm 0$ \citep{Wright-Ma-2022}。\footnote{人们可以通过稍微复杂的分析，直接将我们的结果推广到非正交子空间。}

\begin{figure}[t]
    \begin{subfigure}[t]{0.45\textwidth}
        \centering
        \includegraphics[width=\textwidth]{\toplevelprefix/chapters/chapter5/figs/SNR1.png}
        \caption{噪声水平 $\delta = 0.2$}
    \end{subfigure}
    \hfill
    \begin{subfigure}[t]{0.45\textwidth}
        \centering
        \includegraphics[width=\textwidth]{\toplevelprefix/chapters/chapter5/figs/SNR2.png}
        \caption{噪声水平 $\delta = 0.5$}
    \end{subfigure}
    \caption{{\bf 纯注意力 Transformer 的去噪性能。} 在这里，我们从 \Cref{def:MoG} 中的低秩高斯混合模型中采样初始 token 表示。然后，我们应用 \eqref{eq:MSSA} 来更新 token 表示，并报告每一层的信噪比 (SNR)。}  \label{fig:MSSA}
\end{figure}

现在，令 $\bm Z_k^{\ell}$ 的列表示第 $\ell$ 层来自第 $k$ 个子空间的 token 表示。为了量化去噪能力，我们定义第 $\ell$ 层 token 表示的每个块的信噪比 (SNR) 如下：
\begin{align}\label{def:SNR}
    \mathrm{SNR}(\bm Z_k^{\ell}) \doteq  \frac{\|\bm U_k\bm U_k^T\bm
    Z_k^{\ell} \|_F}{\|(\bm I - \bm U_k\bm U_k^T)\bm Z_k^{\ell}
    \|_F},\quad \forall k \in [K].
\end{align}
为了简化分析，我们假设 $p=p_1=\dots=p_K$，$N_1=\dots=N_K=N/K$，并且
\begin{align}\label{eq:orth}
    \begin{bmatrix}
        \bm U_1 & \dots & \bm U_K
    \end{bmatrix} \in \mathcal{O}^{d\times Kp}.
\end{align}

基于上述设置，我们现在来刻画 MSSA 算子的去噪性能。

\begin{figure*}[t]
    \begin{center}
        \includegraphics[width=0.75\textwidth]{\toplevelprefix/chapters/chapter5/figs/MSSA.png}
        \caption{\textbf{纯注意力 Transformer 架构的详细信息。} 每一层由 MSSA 算子和跳跃连接组成。此外，LayerNorm 仅包含在语言任务中。在实践中，应用反向传播使用训练样本来训练模型参数。 %
        }\label{fig:transformer}
    \end{center}
\end{figure*}

\begin{theorem}\label{thm:1}
    令 $\bm Z^{1}$ 在 \Cref{def:MoG} 中定义，且 \eqref{eq:MSSA} 中的 $\varphi(\cdot)$ 为
    $\varphi(\bm x) = h\left(\sigma(\bm x)\right)$，
    其中 $\sigma:\R^N \to \R^N$ 是 softmax 函数，$h:\R^N
    \to \R^N$ 是逐元素阈值函数，对于每个 $i \in [N]$ 有 $h(x) =
    \tau \mathbb{I}\left\{x > \tau\right\}$。
    假设 $p \gtrsim \log N$，$\delta \lesssim \sqrt{\log
    N}/\sqrt{p}$，并且
    \begin{align*}
        \tau \in \left( \frac{1}{2},  \frac{1}{1+N\exp(-9p/32)} \right].
    \end{align*}
    对于足够大的 $N$，以至少 $1-KLN^{-\Omega(1)}$ 的概率，对于每个 $\ell \in [L]$ 成立：
    \begin{align}\label{eq:SNR}
        \mathrm{SNR}(\bm Z_k^{\ell+1}) = (1+\eta\tau)
        \mathrm{SNR}(\bm Z_k^{\ell}),\ \forall k \in [K].
    \end{align}
\end{theorem}
该定理表明，当初始 token 表示从噪声水平为 $O(\sqrt{\log N}/\sqrt{p})$ 的低秩高斯混合分布中采样时，我们证明了所提出的 Transformer 的每一层都以线性速率对 token 表示进行去噪。这表明了 MSSA 算子在跨层降低噪声方面的高效性。值得注意的是，我们的理论结果得到了 \Cref{fig:MSSA} 中实验观察的有力支持。该定理为通过展开 (\ref{eq:MSSA}) 导出的 Transformer 架构的实际去噪能力提供了理论基础。

\begin{remark}
    在该模型下，表征学习的目标是将一组含噪的初始 token 表示压缩到相应的子空间中。然而，我们需要指出的是，在实际应用中，token 表示表现出更复杂的结构，表征学习的目标是通过压缩 token 集合来找到紧凑且结构化的表示。
\end{remark}

\paragraph{纯注意力 Transformer。} 现在，我们正式提出一种纯注意力 Transformer 架构。具体而言，通过将迭代优化步骤 (\ref{eq:MSSA}) 展开为深度网络的层，我们在 \Cref{fig:transformer} 中构建了一个 Transformer 架构。所提出架构的每一层仅由 MSSA 算子和跳跃连接组成。对于语言任务，我们在 MSSA 算子之前额外加入了 LayerNorm 以提高性能。完整的架构通过堆叠这些层，以及必要的特定任务预处理和后处理步骤（例如位置编码、token 嵌入和适应不同应用的最终特定任务头）来构建。

一般来说，标准的仅解码器 Transformer 架构由以下关键组件构成 \cite{vaswani2017attention}：(1) 位置编码，(2) 多头 QKV 自注意力机制，(3) 前馈 MLP 网络，(4) 层归一化，以及 (5) 残差连接。相比之下，我们提出的 Transformer 架构采用了一种精简设计，包含了几项关键的简化。具体而言，它采用了共享 QKV 子空间自注意力机制，排除了 MLP 层，并降低了 LayerNorm 的使用频率。

\subsection{线性时间注意力：Token 统计 Transformer}\label{sub:tost}

在本小节中，我们基于编码率降低目标提出了一种新的 Transformer 注意力算子，其计算复杂度随 token 数量线性增长。具体而言，我们推导了 MCR$^2$ 目标的一种新颖变分形式，并表明由该变分目标的展开梯度下降产生的架构导致了一个称为 Token 统计自注意力 (\texttt{TSSA}) 的新注意力模块。\texttt{TSSA} 具有{\em 线性的计算和内存复杂度}，并且从根本上背离了计算 token 之间成对相似度的典型注意力架构。回顾 \eqref{eq:MCR pi}，$\bm \Pi = [\bm \pi_1, \ldots, \bm \pi_K] \in \R^{N \times K}$ 表示一个随机的“分组分配”矩阵（即 $\bm \Pi \bm
    1 = \bm 1$ 且 $\Pi_{ik} \geq 0, \  \forall (i,k) \in [N] \times [K]$），其中 $\Pi_{ik}$ 表示将第 $i$ 个 token 分配给第 $k$ 个组的概率。

\paragraph{编码率的新变分形式。} 首先，我们考虑基于矩阵谱的凹函数的类 MCR$^2$ 目标的一般形式。也就是说，对于给定的半正定 (PSD) 矩阵 $\bm M \in \PSD(d)$ 和任意标量 $c \geq 0$，我们有 $\log \det (\I + c \bm M) = \sum_{i=1}^{d} \log( 1 + c \lambda_i(\bm
M))$，其中 $\lambda_i (\bm M)$ 是 $\bm M$ 的第 $i$ 大特征值。此外，请注意 $\log(1 + c \sigma)$ 是关于 $\sigma$ 的凹非递减函数。因此，我们基于形式为 $F(\bm M) = \sum_{i=1}^{d}
f(\lambda_i(\bm M))$ 的 PSD 矩阵的一般谱函数（其中 $f$ 是凹非递减的），用更一般的 MCR$^2$ 形式来描述我们的结果。特别地，回顾我们上面的讨论，注意力机制源于展开 MCR$^2$ 的压缩组件，因此我们考虑一种更一般的 MCR$^2$ 风格的压缩函数：
\vspace{-2mm}
\begin{equation}
    R_{c,f}(\bm Z, \bm \Pi) \doteq \frac{1}{2}\sum_{k=1}^K
    \frac{N_k}{N} F\left( \frac{1}{N_k}\Z \mathrm{Diag}(\bm \pi_k)
    \Z^{\top} \right).
    \label{eq:mcr_c_gen}
\end{equation}

对于上述目标，我们现在注意以下结果：
\begin{theorem}
    \label{thm:var_concave}
    令 \(f \colon [0, \infty) \to \R\) 为非递减、凹函数，且满足 \(f(0) = 0\)，并令 \(F \colon \PSD(d) \to \R\) 具有形式 \(F(\bm M) = \sum_{i = 1}^{d}f(\lambda_{i}(\bm M))\)。
    那么对于每个 \(\bm M \in \PSD(d)\) 和 \(\bm Q \in \O(d)\)，我们有
    \begin{equation}
        \label{eq:U_bound}
        F(\bm M) \leq  \sum_{i=1}^{d} f\left( (\bm Q^{\top} \bm M \bm
        Q)_{ii} \right).
    \end{equation}
    此外，对于任何对角化 $\bm M$ 的 $\bm Q$，\eqref{eq:U_bound} 中的不等式取等号；如果 $f$ 是严格凹的，那么 \eqref{eq:U_bound} 中的不等式取等号\textit{当且仅当} $\bm Q$ 对角化 $\bm M$。
\end{theorem}

利用上述结果，我们可以将 \eqref{eq:mcr_c_gen} 替换为具有以下形式的等效变分目标
\vspace{-2mm}
\begin{equation}
    \label{eq:var_comp}
    R^{\rm var}_{c,f} (\Z,\bm \Pi \mid \bm U_{[K]}) \doteq
    \frac{1}{2}\sum_{k=1}^K \frac{N_k}{N} \sum_{i=1}^{d}
    f\left(\frac{1}{N_k} (\bm U_{k}^{\top} \Z \mathrm{Diag}(\bm
    \pi_k) \bm Z^{\top} \bm U_{k})_{ii} \right),
\end{equation}
这里的等效性是指，对于如 \Cref{thm:var_concave} 中所述的 $\{ \bm U_k \in \O(d)\}_{k=1}^K$ 矩阵的最佳选择（即对角化每个 $\Z \mathrm{Diag}(\bm \pi_k) \Z^\top $ 的正交矩阵），我们将实现一个紧致界 $ R^{\rm var}_{c,f} (\Z,\bm \Pi \mid \bm U_{[K]}) =
R_{c,f} (\Z,\bm \Pi)$。请注意，通常情况下，要实现这个界，需要为每个采样的 $\Z$ 实例选择一组新的最佳 $\bm U_{k}$ 参数矩阵来对角化每个 $\Z\mathrm{Diag}(\bm \pi_{k})\Z^{\top}$，这对于网络架构来说显然是不切实际的。
相反，作为一种替代观点，我们不将数据 ($\Z$) 视为固定的并试图优化 $\bm U_k$ 参数以实现紧致的变分界，而是可以采用上述算法展开设计原则，设计一个算子来扰动 $\Z$，从而增量式地最小化 $R_{c, f}^{\rm
var}(\cdot \mid \bm U_{[K]})$。为了明确这一点，当 $\Z
\mathrm{Diag}(\bm \pi_k) \Z
^\top$ 的特征空间与 $\bm U_k$ 的列对齐时，每个变分界变得紧致，因此通过旋转 $\Z$ 的适当列（即对应于 $\bm \pi_k$ 中较大元素的列）以与 $\bm U_k$ 对齐，我们可以逼近一个紧致的变分界。也就是说，我们不是为每个 $\Z$ 实例旋转 $\bm U_k$ 以与数据对齐，而是可以旋转每个 $\Z$ 中的 token 特征以与 $\bm U_k$ 对齐。

遵循这种方法，我们计算 $R_{c,f}^{\rm var}$ 关于 $\Z$ 的梯度下降步。
为了开始此计算，首先令 $\bm \pi \in \Re^N$ 为任意逐元素非负向量。那么我们有
\begin{equation}\label{eq1:grad Rcf}
    \nabla_{\Z} \ \frac{1}{2} \sum_{i = 1}^{d} f(  (\Z
    \mathrm{Diag}(\bm \pi) \Z^{\top} )_{ii}) %
    = \; \mathrm{Diag}(\nabla f[ \Z^{\hada 2} \bm \pi] ) \Z
    \mathrm{Diag}(\bm \pi),
\end{equation}
其中 $\nabla f$ 是 $f$ 的梯度，并且（回顾）$\nabla
f[\cdot]$ 将 $\nabla f$ 应用于括号中向量的每个元素。特别地，对于 $ f(x) = \log(1 + (d/\epsilon^{2}) x)$，
$\nabla f(x) = (d / \epsilon^{2}) (1+ (d / \epsilon^{2}) x)^{-1}$ 仅仅是一个非线性激活。此外，（回顾）$N_{k} = \langle \bm
\pi_{k}, \bm 1\rangle$。因此，$R^{\rm var}_{c,f}$ 关于 $\Z$ 的梯度为：
\begin{align}\label{eq2:grad Rcf}
    \nabla_{\Z} R^{\rm var}_{c,f}(\Z,\bm \Pi \mid \bm U_{[K]}) =  \frac{1}{n}
    \sum_{k=1}^K \bm U_{k} \underbrace{\mathrm{Diag} \left( \nabla f\left[ (\bm
                U_{k}^{\top} \Z)^{\hada 2}  \frac{\bm \pi_k}{\langle
                    \bm \pi_{k},
    \bm 1\rangle} \right] \right)}_{\doteq \bm D(\Z, \bm \pi_{k} \mid
    \bm U_{k})} \bm U_{k}^{\top} \Z \mathrm{Diag}(\bm \pi_k).
\end{align}
（请注意，$1/N$ 常数源于求和中每一项的 $(N_{k}/N)\cdot (1/N_{k})
= 1/N$ 常数。）如果我们现在考虑关于第 $j$ 个 token $\z_j$ 的梯度步，我们得到了我们提出的增量压缩算子，即我们对\textit{自注意力} + 残差算子的替代：
\vspace{-2mm}
\begin{equation}\label{eq:layer-operation-orig}
    \z_j^+ = \z_{j} - \tau \nabla_{\z_{j}}R_{c, f}^{\rm var}(\Z, \bm
    \Pi \mid \bm U_{[K]}) = \z_j - \frac{\tau}{N} \sum_{k=1}^K
    \Pi_{jk} \bm U_{k} \bm D(\Z, \bm \pi_{k} \mid \bm U_{k}) \bm
    U_{k}^{\top} \z_j
\end{equation}
对于每个 \(j \in [n]\)，其中 $\tau > 0$ 是增量优化的步长参数。然后，我们可以在 \Cref{fig:vcrate-architecture} 中构建一层 TOST。

\begin{figure}[t]
    \centering \includegraphics[width=1\textwidth,trim={0 5.0cm 0
    0},clip]{\toplevelprefix/chapters/chapter5/figs/V-CRATE.pdf}
    \vspace{-1mm}
    \caption{\small \textbf{所提出的 Token 统计 Transformer (ToST) 的第 $\ell$ 层。} 值得注意的是，ToST 的自注意力通过将投影 token 的每一行\textit{仅乘以一个标量}，高效地将 token $\Z^{\ell}$ 转换为 $\Z^{\ell+1}$。这降低了注意力的复杂度：它具有 $O(p)$ 的空间复杂度和 $O(pn)$ 的时间复杂度，其中 $p$ 是每个头投影 token 的维度，$n$ 是 token 的数量。
    }
    \label{fig:vcrate-architecture}
\end{figure}

\paragraph{模型解释。} 给定 \eqref{eq:layer-operation-orig} 中提出的注意力算子，首先回顾 $\bm\Pi$ 的行是非负的且和为 1，因此我们的算子对 $K$ 个类似“注意力头”的算子进行加权平均，然后添加一个残差连接。利用 \(\sum_{k = 1}^{K}\Pi_{jk} = 1\)，我们可以将 \eqref{eq:layer-operation-orig} 重写为： %
\vspace{-2mm}
\begin{equation}
    \label{eq:z_up}
    \z_j^+ = \sum_{k=1}^K \Pi_{jk} \Big[ \z_j \underbrace{-
            \frac{\tau}{n} \bm U_{k} \bm D(\Z, \bm \pi_{k} \mid \bm U_{k})
    \bm U_{k}^{\top}}_{\text{action of one attention head}} \z_j \Big].
\end{equation}
也就是说，我们可以将每个注意力头视为首先通过乘以 $\bm U_k^\top$ 将 token 特征投影到基 $\bm U_{k}$ 上，乘以对角矩阵 $\bm D(\Z, \bm \pi_{k} \mid \bm U_{k})$（简写为 \(\bm
D_{k}\)），通过乘以 $\bm U_{k}$ 投影回标准基，并通过残差连接将其从原始 token 特征中减去。我们注意力层的核心方面是 $\bm D_{k}$ 的计算。即，\(\Pi_{jk} \geq 0\)，因此 $\bm \pi_k / \langle \bm \pi_{k}, \bm
1\rangle \in \Re^N$ 形成了一个关于哪些 token 属于第 $k$ 组的概率分布。因此，$(\bm U^\top_k \Z)^{\hada 2}
\bm \pi_k / \langle \bm \pi_{k}, \bm 1\rangle$ 估计了在由 $\bm
\pi_k /  \langle \bm \pi_{k}, \bm 1\rangle$ 给出的分布下 $\bm U_k^\top \Z$ 的二阶矩。此外，由于 $f$ 是一个凹非递减函数，$\nabla f(x)$ 随着 $x$ 的增加单调递减趋向于 $0$，因此 $\bm D_{k}$ 的元素（具有 $\nabla f(x)$ 的形式）在 $x=0$ 处达到最大值 %
并且随着 $x$ 的增加单调衰减到 $0$。

由此，我们得出了我们的注意力头 + 残差算子 $[\I - (\tau/n)\bm U_{k}\bm D_{k}\bm U_{k}^{\top}]$ 的核心解释。也就是说，该算子执行一种近似的依赖于数据的低秩投影，其中在投影 $\bm
U_k^\top \Z$ 之后具有大量“能量”的方向（即具有较大二阶矩 $(\bm
    U_{k}^{\top}\Z)^{\hada 2}\bm \pi_k / \langle \bm \pi_{k}, \bm
1\rangle$ 的方向）被保留，而不具备这些能量的方向被抑制。为了理解这一点，回顾 $\bm D_k$ 的元素随着二阶矩的增加单调递减到 0，因此对于具有大二阶矩的方向，注意力 + 残差算子在很大程度上充当恒等算子。相反，对于具有小二阶矩的方向，我们的算子减去了 token 沿这些方向的投影，导致这些方向被抑制。与标准的自注意力算子相比，我们的方法显然不计算 token 之间的任何成对相似度。相反，$\Z$ 中 token 之间的交互仅通过它们对用于构建 $\bm D_{k}$ 的二阶矩统计量的贡献来影响算子。尽管如此，与标准的自注意力算子类似，我们的方法仍然具有清晰的解释，即执行一种向依赖于数据的低秩结构压缩的形式，从某种意义上说，它执行了一种近似的低秩投影，其中被抑制的特定方向是那些与组中其他 token 没有强烈对齐的方向。

\paragraph{计算方面的考虑。} 在介绍了我们提出的注意力算子之后，我们现在讨论如何在实践中计算它。首先，在目前的介绍中，我们一直避免讨论如何通过 $\bm\Pi$ 矩阵将 token “分组”到各个注意力头中，但显然需要一种构建 $\bm\Pi$ 的方法来实现我们的方法。此外，我们在 \Cref{thm:var_concave} 中的变分形式要求 $\bm U$ 矩阵是方阵且正交的，但理想情况下，人们希望使用更小的矩阵（即减少 $\bm U$ 中的列数）以提高效率，并放弃正交性约束。

在实践中，我们不强制执行正交性约束。为了减少 $\bm U$ 矩阵中的列数，我们注意到，类似于 CRATE \citep{yu2023white}，如果我们假设组 \(k\) 内的特征 $\bm Z$ （近似）聚集在一个低维子空间周围——假设维度为 \(p\)——那么组 \(k\) 内的协方差 \(\Z\mathrm{Diag}(\bm
\pi_{k})\Z^{\top}\) 是低秩的，回顾 \cite{yu2020learning} 表明 $\Z$ 的最佳几何结构将是每个组都是一个低秩子空间，且与其他组正交。因此，我们可以显式地为该协方差的像（即组 \(k\) 中数据的线性张成）找到一个低维标准正交基。如果维度为 $p \leq d$，该基可以由一个 $d\times p$ 的正交矩阵 $\bm U_k \in \O(d,
p)$ 表示。在这种情况下，我们可以使用这些低秩正交基矩阵更高效地对 \(R_{c,f}\) 给出上界。为了证明这一点，我们使用 \Cref{thm:var_concave} 的更一般版本得出以下推论。
\begin{corollary}\label{cor:var_concave_logdet}
    令 \(f \colon [0, \infty) \to \R\) 为非递减、凹函数，且满足 \(f(0) = 0\)，并令 \(F \colon \PSD(p) \to \R\) 具有形式 \(F(\bm M) = \sum_{i = 1}^{p}f(\lambda_{i}(\bm M))\)。
    令 \(\Z\)、\(\bm \Pi\) 固定。那么，对于所有满足 \(\mathrm{image}(\Z\diag(\bm \pi_{k})\Z^{\top}) \subset
    \mathrm{image}(\bm U_{k})\)（对所有 \(k \in [K]\)）的 \(\bm U_{1},
    \dots, \bm U_{K} \in \O(d, p)\)，我们有
    \begin{equation}
        R_{c, f}(\Z, \bm\Pi) \leq R_{c, f}^{\rm var}(\Z, \bm\Pi \mid
        \bm U_{[K]}),
    \end{equation}
    其中 \(R_{c, f}^{\rm var}\) 在 \eqref{eq:var_comp} 中正式定义。如果 \(\bm U_{k}\) 对角化每个 \(k \in [K]\) 的 \(\Z\diag(\bm \pi_{k})\Z^{\top}\)，则等式成立；如果 \(f\) 是强凹的，那么这个等式条件变为“当且仅当”。
\end{corollary}

定义我们注意力算子的最后一步是估计组成员关系 $\bm\Pi$。为此，我们假设了一个简单的模型，描述每个特征 \(\z_{j}\) 如何偏离其支撑子空间，然后找到最佳的子空间分配。\cite{yu2023white} 表明，如果我们独立地将每个 \(\z_{j}\) 建模为属于一个低维高斯混合模型，其中每个高斯分布具有相同谱的协方差矩阵，并支撑在具有标准正交基 \(\bm U_{k}\) 的子空间上，加上协方差为 \(\eta \I\) 的独立高斯噪声，那么每个 token \(\z_{j}\) 属于每个子空间的后验概率由分配矩阵 \(\bm \Pi = \bm \Pi(\bm Z \mid \bm
U_{[K]})\) 给出如下：
\begin{align}\label{eq:pi}
    \bm \Pi & =
    \begin{bmatrix} \boldsymbol{\nu}(\z_1 \mid \bm U_{[K]})^\top
        \\ \vdots \\ \boldsymbol{\nu}(\z_n \mid \bm U_{[K]})^\top
    \end{bmatrix}, \nonumber \quad
    \text{where} \\
    \boldsymbol{\nu}(\z_j \mid \bm U_{[K]}) &\doteq
    \operatorname{softmax}\left( \frac{1}{2\eta}
        \begin{bmatrix} \|  \bm U_{1}^{\top} \z_j \|_2^2 \\ \vdots
            \\ \| \bm U_{K}^{\top} \z_j \|_2^2
    \end{bmatrix} \right), \quad \forall j \in [n],
\end{align}
其中 $\eta$ 成为一个可学习的温度参数。因此，给定输入特征 \(\Z\)，我们使用 \eqref{eq:pi} 估计 \(\bm\Pi\)，然后计算注意力算子。将 \eqref{eq:pi} 中 $\bm\Pi$ 的构建与 \eqref{eq:layer-operation-orig} 结合起来，我们得到了{\em Token 统计自注意力}算子：
\begin{equation}
    \label{eq:tssa}
    \texttt{TSSA}(\Z \mid \bm U_{[K]}) \doteq -\frac{\tau}{n}\sum_{k
    = 1}^{K}\bm U_{k}\bm D(\Z, \bm\pi_{k} \mid \bm U_{k})\bm
    U_{k}^{\top}\Z\diag(\bm \pi_{k}),
\end{equation}
其中 \(\bm\pi_{k}\) 是在 \eqref{eq:pi} 中定义的 \(\bm\Pi = \bm\Pi(\Z \mid \bm
U_{[K]})\) 的列，\(\bm D\) 在 \eqref{eq2:grad Rcf} 中定义。

\subsection{因果 CRATE}
\label{sec:causal-crate}
深度网络最流行的一些应用是在\textit{序列数据}领域，与前面图像的例子不同，序列数据具有自然的顺序。例如，与图像（人们倾向于一次性查看整张图像）不同，语言是一次处理一个 token（因果地）。这个简单的事实推动了大型语言模型 (LLM) 的发展，这些模型被训练为\textit{仅给定先前 token 的历史}来预测序列中下一个 token 的（概率分布）。然后，为了从 LLM 中采样，只需迭代地从预测的下一个 token 分布中采样，然后将采样的 token 附加到历史记录中。为了训练 LLM 以及其他模型（如视频生成模型），有必要开发能够高效执行这种因果计算的架构。

考虑到这一点，以及从 \Cref{ch:intro} 中对因果自回归过程的理解，我们说一个编码器 \(f\) 是\textit{因果的}，如果对于每个输入 \(\vX = [\vx_{1}, \dots, \vx_{N}]\)，都满足
\begin{equation}
    f(\vX)_{1:N-1} = f([\vx_{1}, \dots, \vx_{N}])_{1:N-1} =
    f([\vx_{1}, \dots, \vx_{N - 1}]) = f(\vX_{1:N-1}),
\end{equation}
其中 \(\vZ_{1:N-1} = [\vz_{1}, \dots, \vz_{N-1}]\) 等等。通俗地说，这意味着前 \(N-1\) 个特征等于编码器在输入的前 \(N-1\) 个 token 嵌入上的输出；更简单地说，这意味着我们可以\textit{一次一个地}计算每个 token 的特征，这等同于计算整个输入的特征。这种因果性对于前面提到的应用（例如 LLM）是必要的。特别是，前面介绍的 CRATE 和 ToST 的实现\textit{不是因果的}（证明留作练习）。

我们可以通过多种方法构建因果白盒架构。在这里，我们将展示一种基于我们之前的展开优化框架的简单方法。具体而言，我们\textit{一次一个地}计算对应于输入前 \(N\) 个 token 的特征 \(\vz_{1}, \dots, \vz_{N}\)，以优化表征学习目标，例如稀疏率降低 \eqref{eq:sparse-rr-1}：
\begin{equation}
    \max_{\vz_{i}}\ \Delta R(\vZ_{1:i} \mid \vU_{[K]}) - \lambda
    \|\vZ_{1:i}\|_{1},\ \forall i \in [N].
\end{equation}
如果我们遵循产生 CRATE 的两步展开过程，我们可以得到以下迭代：
\begin{align}
    \vz_{i}^{\ell + 1/2} &\approx \vz_{i}^{\ell} - \kappa
    \nabla_{\vz_{i}} R(\vZ_{1:i}^{\ell} \mid \vU_{[K]}^{\ell}), \\
    \vz_{i}^{\ell + 1} &\approx \argmin_{\vz_{i}} \left\{\lambda
        \|\vz_{i}\|_{1} + \frac{1}{2}\norm{\vz_{i}^{\ell + 1/2} -
    \vD^{\ell} {\vz_{i}}}_F^2\right\},
\end{align}
或者，使用我们用于 CRATE 的从这两个步骤到网络算子的相同转换，
\begin{align}
    \vz_{i}^{\ell + 1/2}
    &= \left(1 - \frac{\kappa p}{N\epsilon^{2}}\right)\vz_{i}^{\ell}
    + \frac{\kappa p}{N
    \epsilon^{2}}\operatorname{MSSA}(\vZ_{1:i}^{\ell} \mid
    \vU_{[K]}^{\ell})_{i}, \\
    \vz_{i}^{\ell + 1}
    &= \operatorname{ISTA}(\vz_{i}^{\ell + 1/2} \mid \vD^{\ell}).
\end{align}
让我们稍微详细地研究一下这个迭代。首先，请注意，根据构造，这个特征序列对应于一个因果编码器。接下来，假设我们处于这样一种情况：我们已经计算了 \(\vZ_{1:i-1}^{\ell + 1} = [\vz_{1}^{\ell +
1}, \dots, \vz_{i-1}^{\ell + 1}]\)，并在此过程中计算了量
\[(\vU_{k}^{\ell})^{\top}\vZ_{1:i-1}^{\ell} =
    [(\vU_{k}^{\ell})^{\top}\vz_{1}^{\ell}, \dots,
(\vU_{k}^{\ell})^{\top}\vz_{i-1}^{\ell}]\]，并且想要计算 \(\vz_{i}^{\ell + 1}\)（例如与 LLM 推理的情况相关）。在这种情况下，请注意 \(\vz_{i}^{\ell + 1/2}\) 的更新规则可以简化为：
\begin{align}
    \vz_{i}^{\ell + 1/2}
    &= \left(1 - \frac{\kappa p}{N\epsilon^{2}}\right)\vz_{i}^{\ell}
    + \frac{\kappa p}{N
    \epsilon^{2}}\operatorname{MSSA}(\vZ_{1:i}^{\ell} \mid
    \vU_{[K]}^{\ell})_{i} \\
    &= \left(1 - \frac{\kappa p}{N\epsilon^{2}}\right)\vz_{i}^{\ell}
    + \frac{\kappa p^{2}}{N^{2} \epsilon^{4}}\sum_{k =
    1}^{K}\vU_{k}\vU_{k}^{\top}\vZ_{1:i}\softmax((\vU_{k}^{\top}\vZ_{1:i})^{\top}(\vU_{k}^{\top}\vz_{i}))
    \\
    &= \left(1 - \frac{\kappa p}{N\epsilon^{2}}\right)\vz_{i}^{\ell}
    + \frac{\kappa p^{2}}{N^{2} \epsilon^{4}}\sum_{k =
    1}^{K}\vU_{k}\vU_{k}^{\top}[\vZ_{1:i-1},
    \vz_{i}]\softmax((\vU_{k}^{\top}[\vZ_{1:i-1},
    \vz_{i}])^{\top}(\vU_{k}^{\top}\vz_{i})) \\
    &= \left(1 - \frac{\kappa p}{N\epsilon^{2}}\right)\vz_{i}^{\ell}
    + \frac{\kappa p^{2}}{N^{2} \epsilon^{4}}\sum_{k =
    1}^{K}\vU_{k}[\vU_{k}^{\top}\vZ_{1:i-1},
    \vU_{k}^{\top}\vz_{i}]\softmax(([\vU_{k}^{\top}\vZ_{1:i-1},
    \vU_{k}^{\top}\vz_{i}])^{\top}(\vU_{k}^{\top}\vz_{i})).
\end{align}
如果我们缓存前面步骤中的 \(\vU_{k}^{\top}\vZ_{1:i-1}\)，并在每次计算投影 \(\vU_{k}^{\top}\vz_{i}\) 时向其添加单列，那么这一步现在变得\textit{非常高效}。也就是说，我们大大减少了大型矩阵-矩阵乘法的数量，并用缓存加载和矩阵-向量乘法代替它们，这在时间复杂度方面总体上要便宜得多。这种缓存是因果生成模型（如 LLM）能够每秒高效采样数千个 token 的原因，即使每个训练步骤本身需要几秒钟。特征子空间投影的缓存也被称为所谓的“KV 缓存”。

最后，让我们考虑我们要\textit{训练}一个因果 CRATE 模型的情况，并希望在同时给定完整输入序列 \(\vX = [\vx_{1}, \dots, \vx_{N}]\) 的情况下找到最有效的方法。ISTA 步骤可以并行化，成为非因果 CRATE 中的常规 ISTA 步骤，即：
\begin{equation}
    \operatorname{ISTA}(\vZ \mid \vD) \doteq [\ISTA(\vz_{1} \mid
    \vD), \dots, \ISTA(\vz_{N} \mid \vD)],
\end{equation}
因此这意味着 ISTA 步骤与非因果 CRATE 保持一致。MSSA 步骤更有趣，因为它以一种有意义的方式发生了变化。为了了解它是如何变化的，请注意每个 MSSA 算子都有一个值得关注的有趣结构。也就是说，如果我们定义\textit{因果 MSSA 算子}为分块矩阵：
\begin{equation}
    \operatorname{CausalMSSA}(\vZ \mid \vU_{[K]}) \doteq
    [\operatorname{MSSA}(\vZ_{1:1} \mid \vU_{[K]})_{1}, \dots,
    \operatorname{MSSA}(\vZ_{1:N} \mid \vU_{[K]})_{N}],
\end{equation}
然后，通过推导 softmax 代数（证明再次留作练习），我们可以证明
\begin{align}
    \operatorname{CausalMSSA}(\vZ \mid \vU_{[K]})
    &\doteq \frac{p}{N\epsilon^{2}}\mat{\vU_{1}, \dots,
    \vU_{K}}\mat{\softmax((\vU_{1}^{\top}\vZ)^{\top}(\vU_{1}^{\top}\vZ)
        +\vM_{N}) \\ \vdots
    \\ \softmax((\vU_{K}^{\top}\vZ)^{\top}(\vU_{K}^{\top}\vZ) + \vM_{N})}, \\
    \text{where} \quad &\vM_{N} = \mat{0 & -\infty & -\infty & \dots
        & -\infty & -\infty \\ 0 & 0 & -\infty & \dots & -\infty &
        -\infty \\ \vdots & \vdots & \vdots & \ddots & \vdots & \vdots
    \\ 0 & 0 & 0 & \dots & 0 & -\infty \\ 0 & 0 & 0 & \dots & 0 & 0}
    \in \R^{N \times N}.
\end{align}
这里的矩阵 \(\vM_{N}\) 是另一种编码方式，表示 CausalMSSA 中的任何特征都不依赖于未来的特征，因为 softmax 参数的相关项是 \(-\infty\)，因此在 softmax 内的指数运算后被设置为 \(0\)，从而被有效地忽略。这个矩阵 \(\vM_{N}\) 有时被称为\textit{（因果）注意力掩码}。

这样做的结果是，当我们拥有完整序列 \(\vZ^{\ell}\) 时，我们可以在与常规 CRATE 层相似或更短的时间内计算 \(\vZ^{\ell + 1}\)（因为 \(\vM_{N}\) 是硬编码的、与输入无关的，并且使我们能够忽略 softmax 的许多项）。因此，我们可以像常规架构一样定义和训练完整的序列到序列模型。

当我们将此应用于语言建模时（详细信息将在 \Cref{ch:applications} 中提供），我们发现与类似规模的经验设计的语言模型相比，我们可以获得合理的结果，如 \Cref{tab:causal-crate-results} 所示。

\begin{table}[t!]
    \small
    \caption{\small Causal-CRATE-Base 模型与 GPT2-Small、GPT2-Base 模型在数据集测试划分上评估的零样本验证损失（越低越好）。
    }
    \centering
    \resizebox{1\columnwidth}{!}
    {
        \begin{tabular}{ccccccc}
            \hline
            & \#parameters & \textbf{OWT} & \textbf{LAMBADA} &
            \textbf{WikiText} & \textbf{PTB} & \textbf{Avg} \\
            \hline
            GPT2-Base  & {124M} & 2.85 & 4.12 & 3.89 & 4.63 & 3.87 \\
            {GPT2-Small } &  {64M} & {3.04} & {4.49} & {4.31} &
            {5.15} & {4.25} \\
            Causal-CRATE-Base & {60M} & 3.37 & 4.91 & 4.61 & 5.53 & 4.61 \\
            \hline
        \end{tabular}
    }
    \label{tab:causal-crate-results}
\end{table}

尽管因果 CRATE 架构具有理论优势，但在语言建模（以及更一般的序列建模）任务上的差距仍然存在。下面，我们讨论两种可能缩小这一差距的策略；它们各自都是活跃的研究方向。

\begin{example}[更好的目标和展开策略]
    要考虑的一条路径是，稀疏率降低的目标对于因果序列数据可能不是最佳的。例如，我们可能希望确保特征具有某种时间关系或结构，例如，由具有高斯（混合）初始化和高斯步进噪声的线性动力系统生成。最近有一系列爆炸性流行的工作，使用所谓的\textit{状态空间模型}（SSM，其命名与 \Cref{ch:intro} 中讨论的“状态空间模型”的经典概念相悖）\citep{gu2024mamba,yang2024parallelizing,yang2024gated} 高效且显式地对特征中的动力学进行建模。这些模型在经验上很有前景，因为它们避免了对全对注意力的需求，避开了 MSSA 算子及其传统表亲多头注意力的计算和内存复杂度。尽管获得了这种效率提升，但它们具有很高的性能，甚至已被用于一些大规模语言模型中 \citep{team2025kimi}。通过设计更复杂的目标和展开策略，我们也许能够构建一个“白盒” SSM 架构，或 SSM-Transformer 混合架构，从而缩小在因果任务上的差距。
\end{example}

\begin{example}[推理时计算策略]
    改进因果模型的另一条路径是使用更多所谓的\textit{推理时计算}策略：在推理期间（例如从 LLM 采样），找到一种方法来执行更多的某些网络操作，以获得更好的结果。如果我们能够做到这一点，我们就可以通过为手头的任务分配更多的计算资源来解决更难的问题。

    一种朴素的方法是循环模型：将最终特征输入到初始的序列到序列层中，并使其经历更多轮的迭代去噪。这种方法在文献中被称为“循环 Transformer (looped transformers)”（针对主干模型是经验设计的 Transformer 的情况）；这种方法已被证明在一些特别困难的任务中具有优势 \citep{giannou2023looped,yang2023looped,saunshi2025reasoning}。然而，从我们的角度来看，更合理的方法是将迭代在\textit{同一层}中循环多次，在进入下一层之前，针对\textit{同一模型}在每一层采取多个局部去噪或压缩步骤。

    还有一种潜在的方法依赖于我们对特征和模型参数中低维结构的理解。例如，由于 \(\vU_{[K]}\) 矩阵被假定为特征所张成子空间的支撑，我们可以在推理时使用在线算法（如在线 (G)PCA）来更新它们。类似的步骤也许能够在推理时自适应字典。这将使我们能够改进并专门针对特征定制每一层的局部去噪或压缩算子，从而有可能产生更好的性能。

    已经有许多经验性的尝试来改进推理时计算方法。它仍然是一个非常活跃的持续研究领域。
\end{example}

% \subsection{因果 ToST} \label{sec:tost}
% \yima{ToST 的因果版本：与非线性滤波、RNN 和其他线性时间架构的关系...}

\section{总结与注记}

本章介绍的材料基于关于该主题的一系列近期工作，包括 \cite{chan2021redunet,
wang2024global, wang2025attention, wu2025token, yu2023white}。这些贡献涵盖了通过展开优化构建可解释深度网络的理论进展和实用方法。本章讨论的许多关键结果和证明直接源自这些基础性工作，或受其启发。

展开优化算法以构建神经网络的想法可以追溯到开创性工作 \cite{gregor2010learning}。在这项工作中，作者证明了稀疏编码算法——例如迭代收缩阈值算法 (ISTA)——可以被展开以形成多层感知机 (MLP)，从而有效地桥接了迭代优化和神经网络设计。值得注意的是，\cite{monga2019algorithm} 证明了与通用网络相比，这种展开的网络更具可解释性、更高效且更有效。

在本章中，我们基于这一视角，通过展开旨在最小化动机明确的目标（例如前面介绍的（稀疏）率降低目标）的优化算法，来开发有原则的白盒深度网络架构。这种方法不仅阐明了网络中每一层的作用，而且为架构选择提供了理论基础，超越了经验性的试错，走向了可解释和目标驱动的设计。在下文中，我们将通常通过经验设计和启发式调整构建的传统 DNN 与我们具有数学基础的 ReduNet 架构进行比较：

\begin{center}
    \begin{tabular}{| l || c | c |}
        \hline
        & 传统 DNN & ReduNet\\ [0.5ex]
        \hline \hline
        目标 & 输入/输出拟合 & 信息增益\\ [0.5ex]
        \hline
        深度架构 & 试错 & 迭代优化 \\  [0.5ex]
        \hline
        层算子 & 经验性 & 投影梯度 \\  [0.5ex]
        \hline
        平移不变性 & CNN+数据增强 & 不变 ReduNet \\  [0.5ex]
        \hline
        初始化 & 随机/预设计 & 前向展开 \\ [0.5ex]
        \hline
        训练/微调 & 反向传播 & 前向/反向传播\\ [0.5ex]
        \hline
        可解释性 & 黑盒 & 白盒 \\ [0.5ex]
        \hline
        表示 & 隐藏/潜在 & 非相干子空间 \\ [0.5ex]
        \hline
    \end{tabular}
\end{center}

\section{练习与扩展}

\begin{exercise}
    令 $\bm Z = [\bm Z_1,\dots,\bm Z_K] \in \R^{d\times N}$，其中对于每个 $k \in [K]$ 有
    $\bm Z_k \in \R^{d\times N_k}$。对于某个 $\alpha > 0$，令
    \begin{align*}
        R(\bm Z) = \log\det\left(\bm I + \alpha\bm Z\bm Z^\top \right).
    \end{align*}
    1. 给定任意方向 $\bm D \in \R^{d\times m}$，请证明 $\nabla R(\bm Z) = \alpha\bm X^{-1}\bm Z$ 并且
    \begin{align*}
        \nabla^2 R(\bm Z)[\bm D, \bm D] =\alpha \mathrm{Tr}\left( \bm
        X^{-1}\bm D\bm D^\top\right) -
        \frac{\alpha^2}{2}\mathrm{Tr}\left(\bm X^{-1}\left( \bm Z\bm
            D^\top+\bm D\bm Z^\top\right) \bm X^{-1}\left( \bm Z\bm D^\top+\bm
        D\bm Z^\top\right)\right),
    \end{align*}
    其中 $\bm X \doteq \bm I + \alpha\bm Z\bm Z^\top$。{\em 提示：} 注意
    \begin{align*}
        \nabla^2 R(\bm Z)[\bm D, \bm D] \doteq \left\langle \lim_{t
        \to 0} \frac{\nabla R(\bm Z+ t\bm D) - \nabla R(\bm Z)}{t},
        \bm D \right\rangle.
    \end{align*}
    \smallskip

    \noindent
    2. 请证明
    \begin{align*}
        R(\bm Z) \le \sum_{k=1}^K \log\det\left(\bm I + \alpha\bm
        Z_k\bm Z_k^\top \right),
    \end{align*}
    其中等式成立当且仅当对于所有 $k \neq l \in [K]$ 有 $\bm Z_k^\top\bm Z_l = \bm
    0$。
    \bigskip

    \noindent
    3. 给定某个 $\alpha >0$，令对于每个 $k \in [K]$ 有 $\alpha_k=N\alpha/N_k$。请推导以下函数的一阶临界点的闭式解：
    \begin{align*}
        f(\bm Z_k) = \frac{1}{2}\log\det\left(\bm I + \alpha\bm
        Z_k\bm Z_k^\top \right) - \frac{N_k}{2N}\log\det\left(\bm I +
        \alpha_k \bm Z_k\bm Z_k^\top \right) - \frac{\lambda}{2}\|\bm Z_k\|_F^2.
    \end{align*}
    {\em 提示：} 令 $r_k=\mathrm{rank}(\bm Z_k)$。考虑 $\bm Z_k$ 的以下奇异值分解：
    \begin{align*}
        \bm Z_k = \bm P_k\bm \Sigma_k\bm Q_k^\top =
        \begin{bmatrix}
            \bm P_{k,1} & \bm P_{k,2}
        \end{bmatrix}
        \begin{bmatrix}
            \tilde{\bm \Sigma}_{k} & \bm 0 \\
            \bm 0 & \bm 0
        \end{bmatrix}
        \begin{bmatrix}
            \bm Q_{k,1}^\top \\ \bm Q_{k,2}^\top
        \end{bmatrix},
    \end{align*}
    其中 $\bm P_k \in \mathcal{O}^d$ 且 $\bm P_{k,1} \in
    \R^{d\times r_k}$ 和  $\bm P_{k,2} \in \R^{d\times (d-r_k)}$，
    $\bm \Sigma_k \in \R^{d\times N_k}$ 且 $\tilde{\bm \Sigma}_{k}
    \in \R^{r_k\times r_k}$ 为对角矩阵，以及  $\bm Q_k
    \in \mathcal{O}^{N_k}$ 且 $\bm Q_{k,1} \in \R^{N_k\times r_k}$
    和  $\bm P_{k,2} \in \R^{N_k \times (N_k-r_k)}$。
    \medskip

\end{exercise}

\begin{exercise}[矩阵求逆的诺伊曼级数]\label{ex:neumannn}
    令 $\bm A \in \mathbb{R}^{n\times n}$。如果 $\|\bm A\| < 1$，请证明
    \begin{align}\label{eq:neumann}
        \left( \bm I - \bm A\right)^{-1} = \sum_{k=0}^\infty \bm A^k.
    \end{align}
    {\em 提示：} 证明包含两个步骤。 \\
    (i) {\bf 步骤 1}：请使用 $\|\bm A^k\| \le \|\bm A\|^k$ 证明当 $\|\bm A\| < 1$ 时无穷级数
    $\sum_{k=0}^\infty \bm A^k$ 收敛。 \\
    (ii) {\bf 步骤 2}：计算矩阵乘积 $(\bm I - \bm A)
    \sum_{k=0}^\infty \bm A^k$。
\end{exercise}

\begin{exercise}
    请计算 \eqref{eq1:grad Rcf} 和 \eqref{eq2:grad Rcf} 中的梯度。
\end{exercise}

\begin{exercise}
    请证明当 $Kp \le d$ 时的 \Cref{cor:var_concave_logdet}。
\end{exercise}

\begin{exercise}
    推导 TSSA 算子的因果版本。实现它并在因果语言建模任务上将其性能与 CRATE 进行比较。
\end{exercise}

\end{document}